% Options for packages loaded elsewhere
\PassOptionsToPackage{unicode}{hyperref}
\PassOptionsToPackage{hyphens}{url}
\PassOptionsToPackage{dvipsnames,svgnames,x11names}{xcolor}
\documentclass[
  11pt,
  a4paper,
]{article}
\usepackage{xcolor}
\usepackage[margin=25mm]{geometry}
\usepackage{amsmath,amssymb}
\setcounter{secnumdepth}{-\maxdimen} % remove section numbering
\usepackage{iftex}
\ifPDFTeX
  \usepackage[T1]{fontenc}
  \usepackage[utf8]{inputenc}
  \usepackage{textcomp} % provide euro and other symbols
\else % if luatex or xetex
  \usepackage{unicode-math} % this also loads fontspec
  \defaultfontfeatures{Scale=MatchLowercase}
  \defaultfontfeatures[\rmfamily]{Ligatures=TeX,Scale=1}
\fi
\usepackage{lmodern}
\ifPDFTeX\else
  % xetex/luatex font selection
\fi
% Use upquote if available, for straight quotes in verbatim environments
\IfFileExists{upquote.sty}{\usepackage{upquote}}{}
\IfFileExists{microtype.sty}{% use microtype if available
  \usepackage[]{microtype}
  \UseMicrotypeSet[protrusion]{basicmath} % disable protrusion for tt fonts
}{}
\makeatletter
\@ifundefined{KOMAClassName}{% if non-KOMA class
  \IfFileExists{parskip.sty}{%
    \usepackage{parskip}
  }{% else
    \setlength{\parindent}{0pt}
    \setlength{\parskip}{6pt plus 2pt minus 1pt}}
}{% if KOMA class
  \KOMAoptions{parskip=half}}
\makeatother
\usepackage{longtable,booktabs,array}
\newcounter{none} % for unnumbered tables
\usepackage{calc} % for calculating minipage widths
% Correct order of tables after \paragraph or \subparagraph
\usepackage{etoolbox}
\makeatletter
\patchcmd\longtable{\par}{\if@noskipsec\mbox{}\fi\par}{}{}
\makeatother
% Allow footnotes in longtable head/foot
\IfFileExists{footnotehyper.sty}{\usepackage{footnotehyper}}{\usepackage{footnote}}
\makesavenoteenv{longtable}
\setlength{\emergencystretch}{3em} % prevent overfull lines
\providecommand{\tightlist}{%
  \setlength{\itemsep}{0pt}\setlength{\parskip}{0pt}}
\usepackage{microtype}
\usepackage{needspace}
\usepackage{fvextra}
\usepackage{adjustbox}
\AtBeginEnvironment{longtable}{\small\setlength{\tabcolsep}{3pt}}
\DefineVerbatimEnvironment{Highlighting}{Verbatim}{breaklines,commandchars=\\\{\}}
\setlength{\emergencystretch}{3em}
\widowpenalty=10000
\clubpenalty=10000
\allowdisplaybreaks
\usepackage{bookmark}
\IfFileExists{xurl.sty}{\usepackage{xurl}}{} % add URL line breaks if available
\urlstyle{same}
\hypersetup{
  pdftitle={The Grassmannian Degree Law for Passive k-Plane Spectral Routing},
  pdfauthor={Lluis Eriksson (Independent researcher)},
  colorlinks=true,
  linkcolor={blue},
  filecolor={Maroon},
  citecolor={Blue},
  urlcolor={blue},
  pdfcreator={LaTeX via pandoc}}

\title{The Grassmannian Degree Law for Passive k-Plane Spectral Routing}
\author{Lluis Eriksson (Independent researcher)}
\date{9 September 2026 - Internal revision v004}

\begin{document}
\maketitle

\subsection{Abstract}\label{abstract}

A passive lossless \(N\)-port network must send one fixed
\(k\)-dimensional input channel to prescribed output \(k\)-planes
\(Y_1,\allowbreak{}\dots,\allowbreak{}Y_L\) at \(L\) distinct frequencies on the unit circle. We
prove that the minimum McMillan degree \(d_{\min}\) of a rational inner
\(N\times N\) matrix performing this task equals
\(\delta_{\mathrm{Gr}}\), the least degree of a rational curve
\(\mathbb P^1\to\mathrm{Gr}(k,\allowbreak{}H)\) through the data, \(H=\sum_iY_i\),
for every \(k\), \(N\) and \(L\). The lower bound projects the Plücker
vector of the routed column onto \(\Lambda^kH\). The upper bound shows
that the lossless lift \(MH_o^{-1}\) of a Forney minimal basis \(M\) has
exactly \(\sum_j\deg m_j\) states, through the backward-shift identity
\(PK_M=K_{F_0}\), and completes it without new states. The published
laws (\(k=1\) projective degree, direct-sum occupancy \(k(L-n_*)\),
generic \(k(L-1)\)) are special cases. We also prove the additivity
\(\deg(F_0U)=\deg F_0+\deg U\), a base-point law for arbitrary witness
pencils \(P\) (the compiled router has \(\operatorname{fdeg}(P)\) states
plus one per base point in the open disc, so it is minimal iff
\(\operatorname{fdeg}(P)=\delta_{\mathrm{Gr}}\) and no base point lies
in the disc), the generic bound
\(\delta_{\mathrm{Gr}}\ge k\lfloor L(r-k)/r\rfloor\) and, for every
\(L\), the generic value \(k\lfloor L/2\rfloor\) in
\(\mathrm{Gr}(k,\allowbreak{}2k)\), in particular \(2\lfloor L/2\rfloor\) in
\(\mathrm{Gr}(2,\allowbreak{}4)\).

\subsection{1. Introduction}\label{introduction}

\subsubsection{1.1 Antecedents}\label{antecedents}

\textbf{Published special cases.} Three earlier records of the author
determine \(d_{\min}\) on special strata of the data.

\begin{itemize}
\tightlist
\item
  ARR-2026-6M3VGTXZ6W8JW9C9, \emph{Projective Memory and Resonant
  Bottlenecks in Passive Spectral Routing}, Theorem 2.3 (global
  projective memory law): for one input line (\(k=1\)) and target lines
  \(Y_i\) at distinct nodes, \(d_{\min}=\delta_H\), where \(\delta_H\)
  is the least degree of a morphism \(\mathbb P^1\to\mathbb P(H)\) with
  \(R(\zeta_i)=Y_i\), independently of the ambient port number. Its
  lower-bound proof uses the a-priori estimate \(d\le L-1<L\) to force
  the routed column into \(H\).
\item
  ARR-2026-52B6MSS1W197W9T2, \emph{Exact Memory of Finite Spectral
  Routing Tables: The Direct-Sum Occupancy Law and Its Singular Strata},
  Theorem 1.2: for a routing word \(w=w_1\cdots w_L\) over \(m\) target
  \(k\)-planes in direct-sum position (\(\dim\sum_aY_a=mk\)),
  \(d_{\min}=k(L-n_*)\) with \(n_*\) the least occupancy of a letter.
  Its compiler (Lemma 3.1) is only bounded by
  \(k\cdot(\text{column degree})\); Theorem 5.2 gives the detector-rank
  hierarchy \(d_{\min}\ge\max\{r-k,\allowbreak{}\Delta\}\) and Limitation 8.2 states
  that no complete formula is known off the direct-sum locus.
\item
  ARR-2026-1WXCVX96S68GA9VK, \emph{Orthogonal Spectral Fan-Out Costs
  \(k(L-1)\) States}, Theorem 3.1 (mutually orthogonal \(k\)-planes:
  \(d_{\min}=k(L-1)\)), and ARR-2026-33BE6K3HW78F4T7F, \emph{Exact
  Memory of Direct-Sum Spectral Fan-Out}, Theorem 1.1 (universal
  sandwich \(r-k\le d_{\min}\le k(L-1)\), \(r=\dim\sum_iY_i\)) and
  Corollary 1.2 (direct-sum position \(r=kL\) gives \(d_{\min}=k(L-1)\),
  an open dense condition when \(N\ge kL\)).
\end{itemize}

\textbf{The enumerative side.} That the McMillan degree of a linear
system is the degree of a rational curve in a Grassmannian (the graph
curve \(z\mapsto\operatorname{colspan}[F(z);\allowbreak{}I]\)) goes back to
Martin--Hermann {[}MH{]} and is the starting point of
Ravi--Rosenthal--Wang {[}RRW1, RRW2{]} and Sottile {[}So{]}, who
identify the space of degree-\(e\) curves with a Quot scheme and count
curves through Schubert conditions at fixed points of \(\mathbb P^1\) by
Bertram's quantum Schubert calculus {[}Be{]}. On the algebraic side,
Forney {[}Fo{]} showed that a rational vector space has \emph{minimal
polynomial bases} whose column degrees \(e_1\le\dots\le e_k\) are
invariants of the space and that \(\sum_je_j\) is the McMillan degree of
the space (Kailath {[}Ka, §6.5.4{]}).

\textbf{Subspace boundary interpolation.} Boundary Nevanlinna--Pick
interpolation for matrix Schur functions is treated by Bolotnikov--Dym
{[}BD{]}. Rank-one Blaschke--Potapov factorisation {[}Po{]} identifies
the McMillan degree of a square rational inner matrix with the winding
number of its determinant. Alpay, Jorgensen and Lewkowicz {[}AJL,
Theorem 2.1{]} characterise rectangular lossless rational functions by
isometric minimal realizations and complete them to square lossless
functions with the same state space. The matrix Fejér--Riesz theorem is
due to Rosenblum {[}Ro{]}; we use the form in Dritschel--Rovnyak
{[}DR{]}.

The closest antecedent in the filter-design literature is the
\emph{subspace Nevanlinna--Pick interpolation} (SNIP) for discrete-time
matrix all-pass filters of Bharath, Gaharwar, Appaiah and Pal
{[}BGAP{]}. Their boundary problem (their Section 4.1, conditions
(2a)--(2c)) prescribes at \(n\) distinct frequencies \(\omega_i\) the
\emph{full unitary value} \(G(e^{j\omega_i})=A_i\) of an \(m\times m\)
all-pass filter \(G\) together with a positive definite ``group-delay
matrix'' \(\Gamma_i\) unitarily similar to
\(jG^*(e^{j\omega_i})\,\allowbreak{}\tfrac{dG}{d\omega}(e^{j\omega_i})\); the word
``subspace'' refers to the encoding of each value \(A_i\) as the neutral
subspace \(\operatorname{span}\binom{I_m}{A_i}\subset\mathbb C^{2m}\) of
the classical SNIP of control theory, not to subspace-valued data. Their
Theorem 1 states that a rational solution exists iff a modified Pick
matrix with diagonal blocks \(\Gamma_i\) is positive definite, and their
inductive construction (Section 4.3) produces \(G=ND^{-1}\) with \(N,\allowbreak{}D\)
polynomial of degree \(n\), one \(2m\times2m\) linear factor per node,
i.e.~a filter with up to \(nm\) states; the free \(\Gamma_i\) are then
chosen by a semidefinite programme minimising the trace of the Pick
matrix. The routing problem of this paper differs in all three respects:
only the image \(S(\zeta_i)\operatorname{ran}X=Y_i\) of one \(k\)-plane
is prescribed, with bases and phases free; no derivative data enter; and
the question answered is the least McMillan degree compatible with the
data, which {[}BGAP{]} does not pose (its solution has a definite, not a
minimal, degree, and the SNIP data always admit degree-\(nm\) solutions
when the Pick matrix is definite, whereas \(d_{\min}\) here can be as
small as \(r-k\) and as large as \(k(L-1)\)). For interior nodes and
full tangential data, minimal-degree rational interpolation
\emph{without} the inner constraint is classical
(Antoulas--Ball--Kang--Willems {[}ABKW{]}), and inner interpolants of
degree equal to the number of conditions are given by Blaschke--Potapov
products {[}Po{]}; we found no source that states the minimum McMillan
degree of an inner interpolant for subspace-valued boundary data with
free phases (the queries are recorded in the response file of version
2). The comparison with {[}BGAP{]} was made from Sections 1--4 of its
arXiv text.

\subsubsection{1.2 Contribution}\label{contribution}

Theorem 1 states \(d_{\min}=\delta_{\mathrm{Gr}}\) for all \(k\), \(N\)
and \(L\) and all \(k\)-plane data at distinct boundary nodes. The
ingredients that are new relative to the antecedents are:

\begin{enumerate}
\def\labelenumi{\arabic{enumi}.}
\tightlist
\item
  a \emph{degree-preservation lemma} for the lossless lift of a minimal
  basis: \(F_0:=MH_o^{-1}\) has exactly \(\sum_je_j\) states, proved
  through the identity \(PK_M=K_{F_0}\) between backward-shift invariant
  subspaces (Proposition 5.4), which replaces the bound
  \(k\cdot\max_je_j\) of 52B6 Lemma 3.1;
\item
  a \emph{lower bound by coordinate projection of the Plücker vector}
  (Proposition 4.2), valid for every rational \(N\times k\) column
  through the data without any passivity hypothesis, which removes the
  restriction \(d<L\) of 6M3V and shows that \(\delta_{\mathrm{Gr}}\)
  does not depend on \(N\);
\item
  the \emph{additivity} \(\deg(F_0U)=\deg F_0+\deg U\) for inner \(U\)
  (Theorem 5.7), giving uniqueness of the minimum-degree inner column
  with prescribed column space, and the \emph{base-point law} (Theorem
  6.1): the lossless lift of an arbitrary admissible pencil \(P\) has
  \(\operatorname{fdeg}(P)\) states plus one per base point in the open
  disc, and nothing for base points on the circle, outside the closed
  disc, or at infinity; it is a minimum-degree router iff
  \(\operatorname{fdeg}(P)=\delta_{\mathrm{Gr}}\) \emph{and} no base
  point lies in the disc (version 1 omitted the first condition; Example
  6.2);
\item
  the generic lower bound
  \(\delta_{\mathrm{Gr}}\ge k\lfloor L(r-k)/r\rfloor+\max(0,\allowbreak{}k-m_0)\)
  (Proposition 7.3) and the generic value
  \(\delta_{\mathrm{Gr}}=k\lfloor L/2\rfloor\) in \(\mathrm{Gr}(k,\allowbreak{}2k)\)
  for every \(L\) (Theorem 7.4, proved with an explicit witness), in
  particular \(2\lfloor L/2\rfloor\) in \(\mathrm{Gr}(2,\allowbreak{}4)\), which is
  \emph{not} the dimension count \(L-1\); the matching generic upper
  bound for general \((k,\allowbreak{}r)\) is left as a conjecture.
\end{enumerate}

The received version reports AI-assisted drafting and adversarial checks
(see the qualified provenance statement at the end); the present text
includes arguments addressing the uniqueness clause of Lemma C and the
base-point formula. Version 2 answers the external evaluation of the
AIRR workshop of 8 September 2026, which refuted the last clause of
Theorem 6.1 of version 1 by an explicit example (reproduced as Example
6.2), asked that the generic value in \(\mathrm{Gr}(2,\allowbreak{}4)\) be either
restricted to \(3\le L\le7\) or proved for all \(L\) (now proved), and
asked for the exact witnesses behind Table 1 (now in
{\small\texttt{repro/\allowbreak{}witnesses/\allowbreak{}}}); the changes are listed at the end.

\subsection{2. Setting}\label{setting}

\(\mathbb D\) is the open unit disc, \(\mathbb T\) the unit circle,
\(\mathbb C[z]\) the polynomial ring. A rational \(N\times N\) matrix
function \(S\) is \emph{inner} if it is analytic in \(\mathbb D\) and
unitary on \(\mathbb T\). A rational function unitary (or isometric)
almost everywhere on \(\mathbb T\) cannot have a pole on \(\mathbb T\),
so an inner \(S\) is analytic on a neighbourhood of
\(\overline{\mathbb D}\) and contractive there by the maximum principle.

\textbf{McMillan degree.} For a rational \(p\times q\) matrix \(F\)
analytic at \(0\), a \emph{realization} is \(F(z)=D+zC(I-zA)^{-1}B\)
with \(A\in\mathbb C^{n\times n}\); \(\deg F\) is the least such \(n\).
It equals the rank of the block Hankel matrix \((F_{s+t+1})_{s,\allowbreak{}t\ge0}\)
of the Taylor coefficients and the total pole multiplicity of \(F\) on
the Riemann sphere, the point \(\infty\) included; a minimal realization
is reachable and observable, and its \(A\) has
\(\sigma(A)\subset\mathbb D\) when \(F\) is analytic on
\(\overline{\mathbb D}\). We write \(\chi(z):=\det(I-zA)\) for the
realization denominator; \(\chi\ne0\) on \(\overline{\mathbb D}\) in
that case. \(F^\sharp(z):=F(1/\bar z)^*\).

\textbf{Lemma 2.1 (winding).} \emph{If \(U\) is a \(k\times k\) rational
inner function of degree \(n\), then \(\det U\) is a finite Blaschke
product with exactly \(n\) zeros in \(\mathbb D\) (with multiplicity),
i.e.~\(\operatorname{wind}_{\mathbb T}\det U=n\).}

\emph{Proof.} By Potapov's theorem {[}Po{]} a rational inner \(U\) of
degree \(n\) is a product
\(U_0\prod_{q=1}^n\bigl(I+(b_{a_q}-1)u_qu_q^*\bigr)\) of a constant
unitary and \(n\) rank-one factors, \(b_a(z)=(z-a)/(1-\bar az)\),
\(|a_q|<1\), \(\|u_q\|=1\). Each factor has determinant \(b_{a_q}\).
\(\square\) (This is the ``topological McMillan degree'' of
ARR-2026-07EH3CZRD89YC8JN.)

\textbf{Routing data.} An isometric input frame
\(X\in\mathbb C^{N\times k}\); distinct nodes
\(\zeta_1,\allowbreak{}\dots,\allowbreak{}\zeta_L\in\mathbb T\); target \(k\)-planes
\(Y_i\subseteq\mathbb C^N\) (\(\dim Y_i=k\); repeated targets allowed);
\(H:=\sum_iY_i\), \(r:=\dim H\). The \emph{routing problem} asks for
inner \(S\) with
\[S(\zeta_i)\operatorname{ran}X=Y_i\qquad(1\le i\le L),\] bases and
endpoint phases being free, and
\[d_{\min}:=\min\{\deg S:\ S\ \text{rational inner } N\times N,\ S(\zeta_i)\operatorname{ran}X=Y_i\ \forall i\}.\]
The upper bound \(k(L-1)\) of 33BE Theorem 1.1 shows the minimum exists.

\textbf{Chordal error.} For two \(k\)-planes with orthogonal projectors
\(\Pi,\allowbreak{}\Pi'\) we use \(\operatorname{dist}=\|\Pi-\Pi'\|_{\mathrm{op}}\),
the sine of the largest principal angle; for \(k=1\) and unit vectors
this is \(\sqrt{1-|\langle y,\allowbreak{}f\rangle|^2}\), the chordal distance used
in 6M3V. The experiments report
\(\max_i\operatorname{dist}(S(\zeta_i)\operatorname{ran}X,\allowbreak{}Y_i)\).

\subsection{\texorpdfstring{3. Three definitions of
\(\delta_{\mathrm{Gr}}\)}{3. Three definitions of \textbackslash delta\_\{\textbackslash mathrm\{Gr\}\}}}\label{three-definitions-of-delta_mathrmgr}

\textbf{Curves.} A morphism \(R:\mathbb P^1\to\mathrm{Gr}(k,\allowbreak{}H)\)
composed with the Plücker embedding into \(\mathbb P(\Lambda^kH)\) is
given by a vector of \(\binom rk\) homogeneous polynomials of a common
degree \(e\) without common zero on \(\mathbb P^1\); \(e\) is the degree
of \(R\) (the pullback degree of \(\mathcal O(1)\)). Dehomogenising at
\(t=1\) gives a \emph{primitive} polynomial vector
\(v\in\mathbb C[z]\otimes\Lambda^kH\) (coordinates without common zero)
of maximal coordinate degree exactly \(e\), and conversely a primitive
vector of maximal coordinate degree \(e\) whose values are decomposable
defines a morphism of degree \(e\). Put
\[\delta_{\mathrm{Gr}}:=\min\{\deg R:\ R:\mathbb P^1\to\mathrm{Gr}(k,H)\ \text{a morphism},\ R(\zeta_i)=Y_i\ \forall i\}.\]

\textbf{Pencils and minimal bases.} Fix coordinates \(H=\mathbb C^r\).
For a polynomial \(r\times k\) matrix \(P\) of rank \(k\) over
\(\mathbb C(z)\) let
\(\operatorname{fdeg}(P):=\max_I\deg\operatorname{minor}_I(P)-\deg\gcd_I\operatorname{minor}_I(P)\)
(\emph{Forney degree}), the maximum over the \(k\times k\) minors. Let
\(V\subseteq\mathbb C(z)^r\) be a \(k\)-dimensional
\(\mathbb C(z)\)-subspace. A polynomial basis \(M=[m_1\cdots m_k]\) of
\(V\) with column degrees \(e_1\le\dots\le e_k\) is a \emph{minimal
basis} {[}Fo{]} if (a) \(\operatorname{rank}M(z)=k\) for every
\(z\in\mathbb C\) and (b) the leading-coefficient matrix
\([\operatorname{lc}(m_1)\cdots\operatorname{lc}(m_k)]\) has rank \(k\).
Minimal bases exist and the \(e_j\) depend only on \(V\). Finally, for
\(e\ge0\) put
\[V_e:=\{f\in\mathbb C[z]^r:\ \deg f\le e,\ f(\zeta_i)\in Y_i\ \forall i\}.\]

\textbf{Lemma A.}
\[\delta_{\mathrm{Gr}}=\min\{\operatorname{fdeg}(P):\ P\in\mathbb C[z]^{r\times k},\ \operatorname{rank}P(\zeta_i)=k,\ \operatorname{colspan}P(\zeta_i)=Y_i\ \forall i\}\]
\[=\min\Bigl\{\sum_j\deg m_j:\ M\ \text{a minimal basis},\ \operatorname{colspan}M(\zeta_i)=Y_i\ \forall i\Bigr\}\]
\[=\min\Bigl\{\sum_je_j:\ f_j\in V_{e_j},\ f_1(\zeta_i)\wedge\dots\wedge f_k(\zeta_i)\ne0\ \forall i\Bigr\}.\]

\emph{Proof.} (1) \emph{Every polynomial basis \(P\) of \(V\) is
\(P=MU\) with \(U\in\mathbb C[z]^{k\times k}\).} By (a) the
\(k\times k\) minors of \(M\) have no common zero, so a Bézout identity
\(\sum_Ia_I\det M_I=1\) holds with polynomial \(a_I\), and
\(M^+:=\sum_Ia_I\operatorname{adj}(M_I)E_I\) (\(E_I\) the row selector)
is a polynomial left inverse of \(M\); \(U=M^+P\) is polynomial.

\begin{enumerate}
\def\labelenumi{(\arabic{enumi})}
\setcounter{enumi}{1}
\item
  \emph{\(\operatorname{fdeg}(P)=\sum_je_j\) for every polynomial basis
  \(P\) of \(V\).} No minor of \(M\) exceeds degree \(\sum_je_j\), the
  \(z^{\sum e_j}\) coefficient of \(\operatorname{minor}_I(M)\) is the
  corresponding minor of the leading-coefficient matrix, nonzero for
  some \(I\) by (b), and the minors are coprime by (a). By
  Cauchy--Binet,
  \(\operatorname{minor}_I(MU)=\operatorname{minor}_I(M)\det U\), so the
  gcd of the minors of \(P\) is \(\det U\) up to a constant and
  \(\operatorname{fdeg}(P)=\sum_je_j\).
\item
  \emph{Minimal bases are curves and conversely.} The primitive vector
  of minors of \(M\) has maximal degree \(\sum_je_j\) (by (2)), no
  common zero at finite points (by (a)) and none at \(\infty\) (by (b));
  its values are minors of a rank-\(k\) matrix, hence decomposable; so
  it is a morphism of degree \(\sum_je_j\). Conversely, a morphism \(R\)
  of degree \(e\) pulls the tautological subbundle back to a rank-\(k\)
  subbundle \(\mathcal S_R\hookrightarrow\mathcal O^r_{\mathbb P^1}\) of
  degree \(-e\); by Grothendieck's theorem
  \(\mathcal S_R\cong\bigoplus_j\mathcal O(-e_j)\) with \(e_j\ge0\) (an
  injection \(\mathcal O(a)\hookrightarrow\mathcal O\) forces \(a\le0\))
  and \(\sum_je_j=e\). Each inclusion
  \(\mathcal O(-e_j)\to\mathcal O^r\) is a vector of homogeneous
  polynomials of degree \(e_j\); dehomogenising gives \(m_j\) with
  \(\deg m_j\le e_j\), and fibrewise injectivity of the subbundle at
  finite points gives (a), at \(\infty\) that the \(z^{e_j}\)
  coefficients are linearly independent, i.e.~\(\deg m_j=e_j\) and (b).
  The value of \(R\) at \(z\) is \(\operatorname{colspan}M(z)\).
\item
  \emph{Interpolation survives gcd removal.} If \(P=MU\) has
  \(\operatorname{rank}P(\zeta_i)=k\) then \(\det U(\zeta_i)\ne0\), so
  the gcd of the minors of \(P\) does not vanish at the nodes and the
  curve of the primitive vector of minors of \(P\) passes through
  \(\operatorname{colspan}M(\zeta_i)=\operatorname{colspan}P(\zeta_i)=Y_i\).
  In particular \emph{an admissible pencil never has a base point at a
  node.}
\end{enumerate}

Chain: a curve of degree \(\delta_{\mathrm{Gr}}\) gives, by (3), a
minimal basis \(M\) with \(\sum e_j=\delta_{\mathrm{Gr}}\) through the
data, which is an admissible pencil of Forney degree
\(\delta_{\mathrm{Gr}}\) by (2), whose columns lie in \(V_{e_j}\) with
nonzero wedge at every node. Conversely \(f_j\in V_{e_j}\) with nonzero
wedges give a pencil \(P=[f_1\cdots f_k]\) with
\(\operatorname{fdeg}(P)\le\max_I\deg\operatorname{minor}_I(P)\le\sum_je_j\),
and by (4) and (3) the primitive vector of its minors is a curve through
the data of that degree. \(\square\)

The last form is what the exact scripts evaluate (Section 8): for each
tuple \((e_j)\) in increasing \(\sum e_j\), \(V_{e_j}\) is a kernel
computed exactly over \(\mathbb Q(i)\), and for \(k=2\) the existence of
a pair with nonzero wedge at all nodes is equivalent to each of the
\(L\) bilinear forms
\((f_1,\allowbreak{}f_2)\mapsto f_1(\zeta_i)\wedge f_2(\zeta_i)\) on
\(V_{e_1}\times V_{e_2}\) being a nonzero matrix (a finite product of
nonzero polynomials is a nonzero polynomial, so finitely many proper
hypersurfaces do not cover \(V_{e_1}\times V_{e_2}\)).

\textbf{Remark 3.1 (\(N\)-independence).} The minimum over morphisms
\(\mathbb P^1\to\mathrm{Gr}(k,\allowbreak{}\mathbb C^N)\) through the data equals
\(\delta_{\mathrm{Gr}}\): if \(P\) is an \(N\times k\) minimal basis of
degree \(e\) through the data, the \(k\times k\) minors of the
\(r\times k\) block of \(P\) with rows in \(H\) are among the minors of
\(P\), hence of degree \(\le e\), and are nonzero at the nodes because
\(\operatorname{colspan}P(\zeta_i)=Y_i\subseteq H\); dividing by their
gcd gives a curve in \(\mathrm{Gr}(k,\allowbreak{}H)\) of degree \(\le e\) through
the data (decomposability as in Section 4).

\subsection{\texorpdfstring{4. Lower bound:
\(d_{\min}\ge\delta_{\mathrm{Gr}}\)}{4. Lower bound: d\_\{\textbackslash min\}\textbackslash ge\textbackslash delta\_\{\textbackslash mathrm\{Gr\}\}}}\label{lower-bound-d_mingedelta_mathrmgr}

\textbf{Lemma B (minors of a degree-\(n\) function).} \emph{Let
\(F=D+zC(I-zA)^{-1}B\) with \(A\in\mathbb C^{n\times n}\) and
\(\chi=\det(I-zA)\). For every \(j\times j\) submatrix \(F_{IJ}\),
\(\chi\det F_{IJ}\) is a polynomial of degree \(\le n\).}

\emph{Proof.} The Schur-complement identity
\[\det\begin{pmatrix}I-zA&-zB_J\\ C_I&D_{IJ}\end{pmatrix}=\det(I-zA)\,\det\bigl(D_{IJ}+zC_I(I-zA)^{-1}B_J\bigr)=\chi\det F_{IJ}\]
exhibits \(\chi\det F_{IJ}\) as the determinant of an
\((n+j)\times(n+j)\) matrix whose first \(n\) rows are affine in \(z\)
and whose last \(j\) rows are constant; by multilinearity in the rows it
is a polynomial of degree \(\le n\). \(\square\)

No inner-ness is used; 07EH Lemma 3.2, 1WXC Lemma 3.2 and 52B6 Lemma 2.1
are this statement for inner \(F\). Poles at \(\infty\) (singular \(A\))
are included in the count \(n\).

\textbf{Proposition 4.2 (projection bound).} \emph{Let \(F\) be a
rational \(N\times k\) matrix function analytic at \(0\) and at the
nodes with \(\operatorname{colspan}F(\zeta_i)=Y_i\) for all \(i\). Then
\(\deg F\ge\delta_{\mathrm{Gr}}\).}

\emph{Proof.} Let \(n=\deg F\) and take a minimal realization, so
\(\chi(\zeta_i)\ne0\) (the poles of \(F\) are the zeros of \(\chi\)
together with possibly \(\infty\)). Choose a unitary \(U\) with
\(U^*H=\mathbb C^r\oplus0\) and replace \((F,\allowbreak{}Y_i)\) by
\((U^*F,\allowbreak{}U^*Y_i)\); so assume \(H=\operatorname{span}(e_1,\allowbreak{}\dots,\allowbreak{}e_r)\).
Let
\[w(z):=\chi(z)\cdot\bigl(\text{the }k\times k\text{ minors of }F(z)\text{ with rows in }\{1,\dots,r\}\bigr)\in\mathbb C[z]\otimes\Lambda^k\mathbb C^r .\]
By Lemma B, \(w\) is a polynomial vector of degree \(\le n\); it is
\(\chi\) times the Plücker vector of the top \(r\times k\) block
\(F_H(z)\) of \(F(z)\). At a node,
\(\operatorname{colspan}F(\zeta_i)=Y_i\subseteq H\) forces the last
\(N-r\) rows of \(F(\zeta_i)\) to vanish, so \(F_H(\zeta_i)\) has rank
\(k\) and
\(w(\zeta_i)=\chi(\zeta_i)\cdot(\text{nonzero multiple of the Plücker vector of }Y_i)\ne0\).
Hence \(w\not\equiv0\); let \(g\) be the gcd of its coordinates and
\(v:=w/g\). Then \(v\) is primitive of degree \(\le n\);
\(g(\zeta_i)\ne0\), so \(v(\zeta_i)\) is a nonzero multiple of the
Plücker vector of \(Y_i\); and for every \(z\) with \(\chi(z)g(z)\ne0\),
\(v(z)\) is a scalar multiple of the vector of minors of \(F_H(z)\),
hence lies in the cone over \(\mathrm{Gr}(k,\allowbreak{}r)\) (decomposable vectors,
including \(0\)), which is Zariski-closed in \(\Lambda^k\mathbb C^r\). A
polynomial map whose values lie in a closed set off a finite set lies in
it everywhere, also at \(\infty\) after homogenising. So \(v\) is a
morphism \(\mathbb P^1\to\mathrm{Gr}(k,\allowbreak{}H)\) of degree \(\le n\) through
the data, and \(\delta_{\mathrm{Gr}}\le n\). \(\square\)

\textbf{Corollary 4.3.} \(d_{\min}\ge\delta_{\mathrm{Gr}}\).

\emph{Proof.} If \(S\) is an inner interpolant of degree \(d\) with
realization \(D+zC(I-zA)^{-1}B\), then \(F:=SX=DX+zC(I-zA)^{-1}(BX)\)
has a realization with \(d\) states, so \(\deg F\le d\); \(F\) is
analytic on \(\overline{\mathbb D}\) and
\(\operatorname{colspan}F(\zeta_i)=S(\zeta_i)\operatorname{ran}X=Y_i\).
Apply Proposition 4.2. \(\square\)

\textbf{Remarks 4.4.} (i) The relation between \(d\) and \(L\) plays no
role. The 6M3V argument (``\(\ell n\) vanishes at \(L>d\) nodes, so the
column stays in \(H\)'') fails for \(k>1\), where \(d_{\min}\) can
exceed \(L\) (e.g.~\(k(L-1)\)); it is not needed, because the column may
leave \(H\) along the curve of \(F\) while its Plücker projection to
\(\Lambda^kH\) still has degree \(\le n\), and removing the gcd can only
lower the degree. (ii) The bound holds for every rational column through
the data, passive or not; passivity enters only in the upper bound, and
the proposition applies verbatim to the rectangular inner column \(F_0\)
of Section 5.

\subsection{\texorpdfstring{5. Upper bound:
\(d_{\min}\le\delta_{\mathrm{Gr}}\)}{5. Upper bound: d\_\{\textbackslash min\}\textbackslash le\textbackslash delta\_\{\textbackslash mathrm\{Gr\}\}}}\label{upper-bound-d_minledelta_mathrmgr}

Throughout, \(M=[m_1\cdots m_k]\) is a minimal basis in
\(H\)-coordinates (\(r\times k\)) with
\(\operatorname{colspan}M(\zeta_i)=Y_i\) and
\(\sum_je_j=\delta_{\mathrm{Gr}}\) (Lemma A), and \(m:=\max_je_j\).

\textbf{Theorem 5.1 (matrix Fejér--Riesz; Rosenblum {[}Ro{]},
Dritschel--Rovnyak {[}DR{]}).} \emph{Let \(R(z)=\sum_{|t|\le m}R_tz^t\)
be a \(k\times k\) Laurent polynomial with \(R(z)\succeq0\) on
\(\mathbb T\) and \(\det R\not\equiv0\). Then \(R=H^\sharp H\) on
\(\mathbb T\) for a polynomial \(H(z)=\sum_{t=0}^mH_tz^t\) that is
outer: \(\det H(z)\ne0\) for \(z\in\mathbb D\). If \(R\succ0\) on
\(\mathbb T\) then also \(\det H\ne0\) on \(\mathbb T\).}

The last clause follows from \(|\det H|^2=\det R\) on \(\mathbb T\).
Apply it to \(R:=M^\sharp M\), \(R_t=\sum_sM_s^*M_{s+t}\), which is
\(\succ0\) on \(\mathbb T\) because \(M(z)\) has rank \(k\) at
\emph{every} \(z\in\mathbb C\) (property (a) of a minimal basis, not
only at the nodes). Let \(H_o\) be the outer factor and put
\[F_0:=MH_o^{-1}.\] \(F_0\) is rational and analytic on a neighbourhood
of \(\overline{\mathbb D}\), \(F_0^*F_0=H_o^{-*}M^*MH_o^{-1}=I_k\) on
\(\mathbb T\) (a rectangular inner column), and
\(\operatorname{colspan}F_0(\zeta_i)=\operatorname{colspan}M(\zeta_i)=Y_i\)
since \(H_o(\zeta_i)\) is invertible.

\textbf{Backward shift.} On the Hardy space \(H^2_N\) let
\((S^*h)(z)=(h(z)-h(0))/z\). For a rational
\(F\in H^\infty_{N\times k}\) put
\[K_F:=\operatorname{span}\{S^{*m}(Fc):\ m\ge1,\ c\in\mathbb C^k\}\subset H^2_N .\]

\textbf{Lemma 5.2.} \emph{(i) \(\dim K_F=\deg F\). (ii) For a polynomial
matrix \(M\) with column degrees \(e_j\), \(\dim K_M\le\sum_je_j\).}

\emph{Proof.} (i) With a minimal \(n\)-state realization,
\(F=D+\sum_{m\ge1}z^mCA^{m-1}B\), so
\(S^{*m}F=\sum_{t\ge0}z^tCA^{m+t-1}B=C(I-zA)^{-1}A^{m-1}B\). Hence
\(K_F=\{C(I-zA)^{-1}x:\ x\in\operatorname{span}_{m,\allowbreak{}c}A^{m-1}Bc\}=\{C(I-zA)^{-1}x:\ x\in\mathbb C^n\}\)
by reachability, and \(x\mapsto C(I-zA)^{-1}x=\sum_tz^tCA^tx\) is
injective by observability. (ii) \(S^{*m}(Me_j)=0\) for \(m>e_j\), so
the \(\sum_je_j\) vectors \(S^{*m}(Me_j)\), \(1\le m\le e_j\), span
\(K_M\). \(\square\)

\textbf{Lemma 5.3 (product rule).} \emph{For
\(f\in H^\infty_{N\times k}\) and \(g\in H^2_k\),
\(S^*(fg)=f\cdot S^*g+(S^*f)g(0)\), and for \(h=\sum_th_tz^t\in H^2_k\)
and \(m\ge1\)}
\[S^{*m}(fh)=f\cdot S^{*m}h+\sum_{j=1}^m(S^{*j}f)\,h_{m-j}.\]

\emph{Proof.} The first identity is
\(\frac{fg-f(0)g(0)}z=f\frac{g-g(0)}z+\frac{f-f(0)}zg(0)\). Induction on
\(m\): applying \(S^*\) to the displayed formula,
\(S^*(fS^{*m}h)=fS^{*(m+1)}h+(S^*f)(S^{*m}h)(0)\) with
\((S^{*m}h)(0)=h_m\), and
\(S^*\bigl((S^{*j}f)h_{m-j}\bigr)=(S^{*(j+1)}f)h_{m-j}\) because
\(h_{m-j}\) is constant. \(\square\)

\textbf{Proposition 5.4 (\(PK_M=K_{F_0}\)).} \emph{Let \(P\) be the
orthogonal projection of \(H^2_N\) onto \((F_0H^2_k)^\perp\). Then
\(PK_M=K_{F_0}\).}

\emph{Proof.} \emph{(i) \(K_{F_0}\perp F_0H^2_k\).} Multiplication by
\(F_0\) is an isometry \(H^2_k\to H^2_N\) since \(F_0^*F_0=I\) on
\(\mathbb T\). For \(m\ge1\), \(c\in\mathbb C^k\), \(h\in H^2_k\):
\(\langle S^{*m}(F_0c),\allowbreak{}F_0h\rangle=\langle F_0c,\allowbreak{}z^mF_0h\rangle=\langle c,\allowbreak{}z^mh\rangle=0\).
Hence \(P|_{K_{F_0}}=\mathrm{id}\).

\emph{(ii) \(PK_M\subseteq K_{F_0}\).} For \(c\in\mathbb C^k\) put
\(h:=H_oc\), so \(F_0h=Mc\). Lemma 5.3 gives
\(S^{*m}(Mc)=F_0S^{*m}h+\sum_{j=1}^m(S^{*j}F_0)h_{m-j}\); \(P\) kills
the first term, so
\(PS^{*m}(Mc)=\sum_{j=1}^m(S^{*j}F_0)(H_oc)_{m-j}\in K_{F_0}\).

\emph{(iii) \(K_{F_0}\subseteq PK_M\).} For \(\lambda\in\mathbb D\) and
\(c\in\mathbb C^k\) set
\[\begin{aligned}\Phi_c(\lambda)&:=\sum_{m\ge1}\lambda^mPS^{*m}(Mc)\\&=\sum_{j\ge1}\lambda^j(S^{*j}F_0)\sum_{t\ge0}\lambda^t(H_o)_tc=G(\lambda)H_o(\lambda)c,\\G(\lambda)&:=\sum_{j\ge1}\lambda^jS^{*j}F_0.\end{aligned}\]
The sum defining \(\Phi_c\) is finite (\(S^{*m}(Mc)=0\) for
\(m>m_{\max}\)); \(G(\lambda)=C(I-zA)^{-1}\lambda(I-\lambda A)^{-1}B\)
converges for \(|\lambda|<1/\rho(A)\), in particular on \(\mathbb D\),
and takes values in the finite-dimensional space \(K_{F_0}\), so the
rearrangement is legitimate. For an analytic map \(\Psi\) from
\(\mathbb D\) into a finite-dimensional space, the span of its values
equals the span of its Taylor coefficients at \(0\) (a linear functional
annihilating one set annihilates the other, and subspaces of a
finite-dimensional space are determined by their annihilators).
Therefore
\[\begin{aligned}PK_M&=\operatorname{span}_{m,c}PS^{*m}(Mc)=\operatorname{span}_{\lambda,c}\Phi_c(\lambda)\\&=\operatorname{span}_{\lambda,c}G(\lambda)H_o(\lambda)c=\operatorname{span}_{\lambda,c'}G(\lambda)c'\\&=\operatorname{span}_{j\ge1,c'}(S^{*j}F_0)c'=K_{F_0}.\end{aligned}\]
using that \(H_o(\lambda)\) is invertible for every
\(\lambda\in\mathbb D\) (outerness) in the fourth equality and the same
span principle for \(G\) in the fifth. \(\square\)

\textbf{Corollary 5.5.}
\(\deg F_0=\dim K_{F_0}=\dim PK_M\le\dim K_M\le\sum_je_j=\delta_{\mathrm{Gr}}\);
combined with Proposition 4.2, \(\deg F_0=\delta_{\mathrm{Gr}}\)
exactly, and all inequalities are equalities (\(\dim K_M=\sum_je_j\),
\(P\) is injective on \(K_M\)).

The degree identity also holds for an arbitrary minimal basis, without
assuming that it minimizes the interpolation problem: if \(M\) has total
column degree \(e\), Proposition 5.4 gives \(\deg(MH_o^{-1})\le e\).
Conversely, take a minimal realization of \(F_0=MH_o^{-1}\) of degree
\(n\) and denominator \(\chi\). Lemma B makes the vector
\(w_I=\chi\operatorname{minor}_I(F_0)\) polynomial of degree at most
\(n\). Its primitive vector represents the same rational column space as
\(M\), so Lemma A(2) gives its degree as \(e\). Thus \(e\le n\), and
\(\deg F_0=e\) for every minimal basis. This is the version used in
Theorem 5.7, where the rational column space is fixed independently of
the interpolation minimum.

\textbf{Theorem 5.6 (same-state completion; Alpay--Jorgensen--Lewkowicz
{[}AJL, Theorem 2.1 I(i)--(ii), p.~5, equations (2.3)--(2.5){]}).}
\emph{Let \(F\) be a \(p\times m\) rational function, \(p\ge m\), with
all poles in \(\mathbb D\) (Schur stable), and write
\(F(w)=D+C(wI-A)^{-1}B\). (i) \(F\) is isometric on \(\mathbb T\) iff it
admits a minimal realization matrix
\(R=\begin{pmatrix}A&B\\ C&D\end{pmatrix}\) with \(R^*R=I_{n+m}\). (ii)
If so, there exist \(\tilde B\in\mathbb C^{n\times(p-m)}\),
\(\tilde D\in\mathbb C^{p\times(p-m)}\) such that
\(\begin{pmatrix}A&B&\tilde B\\ C&D&\tilde D\end{pmatrix}\) is unitary.}

Part (ii) is elementary: an isometry is completed to a unitary by an
orthonormal basis of the orthogonal complement of its range; this is
what the numerical pipeline does. To apply the theorem, note that if
\(F_0(z)=D+zC(I-zA)^{-1}B\) is our disc-convention realization, then
\(F_0(1/w)=D+C(wI-A)^{-1}B\) has the \emph{same} realization matrix, is
analytic at \(w=\infty\), has all poles in \(\mathbb D\) (\(F_0\) is
analytic on \(\overline{\mathbb D}\)), is isometric on \(\mathbb T\)
(\(\mathbb T\) is invariant under \(w\mapsto1/w\)) and is minimal iff
\(F_0\) is. No paraconjugation is needed. In AJL's notation \(p=r\),
\(m=k\) and the state dimension is \(n\); equation (2.3) makes the
\((n+r)\times(n+k)\) matrix an isometry, and (2.4)--(2.5) add exactly
\(r-k\) input columns without changing \(A\) or \(C\). This identifies
the same-state assertion used below.

\textbf{Theorem 1 (Grassmannian degree law).} \emph{For every \(k\),
every \(N\), every finite table of \(k\)-planes at distinct boundary
nodes,} \[d_{\min}=\delta_{\mathrm{Gr}} .\] \emph{The value depends only
on the configuration \((Y_i)\subset\mathrm{Gr}(k,\allowbreak{}H)\) and the nodes, not
on \(N\). A minimum-degree curve, represented by a minimal basis \(M\),
compiles to a square inner router of the same degree by
\(M\mapsto MH_o^{-1}\) followed by unitary augmentation of an isometric
realization; conversely every router of degree \(d\) yields a curve of
degree \(\le d\) through the data.}

\emph{Proof.} Corollary 4.3 gives \(\ge\). For \(\le\), let
\(n=\deg F_0=\delta_{\mathrm{Gr}}\) (Corollary 5.5) and take, by Theorem
5.6 (i)--(ii), a minimal isometric realization of \(F_0\) augmented to a
unitary matrix. Then
\(S_H(z):=[D\ \tilde D]+zC(I-zA)^{-1}[B\ \tilde B]\) is \(r\times r\)
with a unitary realization matrix and \(\sigma(A)\subset\mathbb D\),
hence analytic on \(\overline{\mathbb D}\) and unitary on \(\mathbb T\)
by the identity
\(I-S_H(z)^*S_H(z)=(1-|z|^2)B_f^*(I-\bar zA^*)^{-1}(I-zA)^{-1}B_f\),
\(B_f=[B\ \tilde B]\); it has \(\le n\) states, and its first \(k\)
columns are \(F_0\) of degree \(n\), so \(\deg S_H=n\). Let \(U\) be
unitary with \(U(\mathbb C^r\oplus0)=H\) and \(V_0\) unitary with
\(V_0\operatorname{ran}X=\mathbb C^k\oplus0\), and set
\[S:=U\,(S_H\oplus I_{N-r})\,V_0 .\] Then
\(\deg S=n=\delta_{\mathrm{Gr}}\) and
\(S(\zeta_i)\operatorname{ran}X=U\bigl(\operatorname{colspan}F_0(\zeta_i)\oplus0\bigr)=Y_i\).
\(\square\)

\textbf{Theorem 5.7 (Lemma C: additivity and uniqueness).} \emph{Let
\(F\) be an inner \(N\times k\) column (analytic on
\(\overline{\mathbb D}\), isometric on \(\mathbb T\)) whose
\(\mathbb C(z)\)-column space is \(V\), let \(M\) be a minimal basis of
\(V\) with column degrees \(e_j\) and \(F_0=MH_o^{-1}\). Then \(F=F_0U\)
with \(U\) a \(k\times k\) rational inner function and}
\[\deg F=\deg F_0+\deg U=\sum_je_j+\deg U .\] \emph{In particular
\(F_0\) is the unique minimum-degree inner column with column space
\(V\), up to a constant unitary factor on the right.}

\emph{Proof.} Write \(F=P/\chi\) with \(\chi\) the denominator of a
minimal realization (\(\deg\chi\le n':=\deg F\), \(\chi\ne0\) on
\(\overline{\mathbb D}\)) and \(P\) polynomial with column space \(V\)
(over \(\mathbb C(z)\), \(F\) has rank \(k\) since it is isometric on
\(\mathbb T\)). By Lemma A(1), \(P=MQ\) with \(Q\) polynomial, so
\(F=MQ/\chi=F_0\,\allowbreak{}(H_oQ/\chi)=:F_0U\). \(U\) is rational and analytic on
\(\overline{\mathbb D}\), and on \(\mathbb T\),
\(U^*U=U^*F_0^*F_0U=F^*F=I\); so \(U\) is inner. The inequality
\(\deg(F_0U)\le\deg F_0+\deg U\) is the product of realizations. For the
reverse, Cauchy--Binet gives
\(\operatorname{minor}_I(F)=\operatorname{minor}_I(M)\det Q/\chi^k\),
and Lemma B applied to the minimal realization of \(F\) says that
\(\chi\operatorname{minor}_I(F)=\operatorname{minor}_I(M)\det Q/\chi^{k-1}\)
is a polynomial of degree \(\le n'\) for every \(I\). Thus
\(\chi^{k-1}\) divides \(\operatorname{minor}_I(M)\det Q\) for all
\(I\); since the minors of \(M\) are coprime (Bézout),
\(\chi^{k-1}\mid\det Q\), say \(\det Q=\chi^{k-1}\tilde q\), and
\(\deg\operatorname{minor}_I(M)+\deg\tilde q\le n'\) for all \(I\).
Taking \(I\) with \(\deg\operatorname{minor}_I(M)=\sum_je_j\) (Lemma
A(2)) gives \(n'\ge\sum_je_j+\deg\tilde q\). Now
\(\det U=\det H_o\det Q/\chi^k=\det H_o\,\allowbreak{}\tilde q/\chi\), and by Lemma
2.1 \(\det U\) has exactly \(\deg U\) zeros in \(\mathbb D\); none can
come from \(\det H_o\) (outer, nonzero on \(\mathbb T\)) or from
\(1/\chi\) (zero-free on \(\overline{\mathbb D}\)), so
\(\deg\tilde q\ge\#\{\text{zeros of }\tilde q\text{ in }\mathbb D\}\ge\deg U\).
Hence \(n'\ge\sum_je_j+\deg U=\deg F_0+\deg U\). If \(\deg F=\deg F_0\)
then \(\deg U=0\): \(U\) is a constant unitary. \(\square\)

\subsection{6. Base points}\label{base-points}

Theorem 1 uses a minimal basis, which has no base points. In practice a
witness pencil \(P\) may be produced by any means (Lagrange
interpolation, a random element of
\(V_{e_1}\times\dots\times V_{e_k}\)), and the natural compiler is
\(P\mapsto PH_P^{-1}\) with \(H_P\) the outer factor of \(P^\sharp P\)
(Theorem 5.1 in its semidefinite form; \(H_P\) may then be singular at
points of \(\mathbb T\)). The following statement, asserted by
experiment in the first version and proved here, explains what the
compiler returns.

\emph{Boundary convention.} All rational products \(PH_P^{-1}\) below
are understood after cancelling removable singularities. The outer
factor has no zero in \(\mathbb D\). If \(\det H_P\) vanishes at
\(\xi\in\mathbb T\), then on a punctured boundary arc the identity
\((PH_P^{-1})^*(PH_P^{-1})=I_k\) bounds every entry by one. An entry
with a genuine pole at \(\xi\) would be unbounded along that arc. Thus
each possible pole is removable, and the isometry identity extends by
continuity at \(\xi\). There is no extra McMillan state from such a
cancelled boundary factor.

\textbf{Theorem 6.1 (base-point law).} \emph{Let
\(P\in\mathbb C[z]^{r\times k}\) be an admissible pencil
(\(\operatorname{rank}P(\zeta_i)=k\),
\(\operatorname{colspan}P(\zeta_i)=Y_i\)), \(V\) its column space, \(M\)
a minimal basis of \(V\), \(F_0=MH_o^{-1}\), and \(g_P\) the gcd of the
\(k\times k\) minors of \(P\) (its zeros are the base points of \(P\);
none lies at a node by Lemma A(4)). Then \(PH_P^{-1}=F_0U_i\) with
\(U_i\) a \(k\times k\) rational inner function, and}
\[\deg(PH_P^{-1})=\operatorname{fdeg}(P)+\#\{\text{zeros of }g_P\text{ in }\mathbb D\}\quad(\text{with multiplicity}).\]
\emph{Base points on \(\mathbb T\setminus\{\zeta_i\}\), outside
\(\overline{\mathbb D}\), or at \(\infty\) (a drop of rank of the
leading-coefficient matrix) are free; base points in \(\mathbb D\) cost
one state each. Since \(\operatorname{fdeg}(P)\ge\delta_{\mathrm{Gr}}\)
for every admissible \(P\) (Lemma A), the compiler attains \(d_{\min}\)
iff \(\operatorname{fdeg}(P)=\delta_{\mathrm{Gr}}\) and \(P\) has no
base point in \(\mathbb D\).} The first condition holds automatically
when \(P=[f_1\cdots f_k]\) with \(f_j\in V_{e_j}\),
\(\sum_je_j=\delta_{\mathrm{Gr}}\) and nonzero wedges at the nodes (the
witnesses of Section 8; Lemma A gives
\(\operatorname{fdeg}(P)\le\sum_je_j\)), but not for an arbitrary
admissible pencil, which merely interpolates (Example 6.2).

\emph{Proof.} By Lemma A(1), \(P=MU_P\) with \(U_P\) polynomial,
\(\det U_P=g_P\) up to a constant. Then
\(P^\sharp P=U_P^\sharp H_o^\sharp H_oU_P=(H_oU_P)^\sharp(H_oU_P)\) on
\(\mathbb T\). Let \(H_oU_P=U_iH_P'\) be the inner--outer factorisation
of the polynomial matrix \(H_oU_P\) (\(\det\not\equiv0\)): \(U_i\) is a
square rational inner function and \(H_P'\) is outer. Then
\(P^\sharp P=H_P'^\sharp H_P'\) on \(\mathbb T\), so \(H_P'\) is an
outer spectral factor of \(P^\sharp P\) and coincides with \(H_P\) up to
a constant unitary on the left; absorbing it into \(U_i\),
\[PH_P^{-1}=MU_PH_P^{-1}=MH_o^{-1}(H_oU_P)H_P^{-1}=F_0U_i .\] By Theorem
5.7,
\(\deg(PH_P^{-1})=\deg F_0+\deg U_i=\operatorname{fdeg}(P)+\deg U_i\)
(Lemma A(2)). By Lemma 2.1, \(\deg U_i\) is the number of zeros of
\(\det U_i=\det H_o\det U_P/\det H_P\) in \(\mathbb D\); \(\det H_o\)
and \(\det H_P\) have no zeros in \(\mathbb D\) (outer), so this is the
number of zeros of \(\det U_P=g_P\) in \(\mathbb D\). Zeros of \(g_P\)
on \(\mathbb T\) are absorbed by the outer factor (\(\det H_P\) may
vanish on \(\mathbb T\)) and do not count; zeros outside
\(\overline{\mathbb D}\) and a rank drop at \(\infty\) do not contribute
an additional term to the degree formula beyond
\(\operatorname{fdeg}(P)\). For the final claim, Lemma A says that
\(\delta_{\mathrm{Gr}}\) is the \emph{minimum} of
\(\operatorname{fdeg}\) over admissible pencils, so
\(\operatorname{fdeg}(P)\ge\delta_{\mathrm{Gr}}\) and
\(\deg(PH_P^{-1})\ge\delta_{\mathrm{Gr}}=d_{\min}\) (Theorem 1), with
equality iff both terms of the formula are minimal. (Version 1 asserted
at this point that every admissible \(P\) has
\(\operatorname{fdeg}(P)=\delta_{\mathrm{Gr}}\), and concluded that the
compiler is minimal iff there is no base point in \(\mathbb D\); the
workshop example below refutes that clause, while the degree formula
stands.) \(\square\)

\textbf{Example 6.2 (an admissible pencil without base points whose lift
is not minimal; AIRR workshop, 2026-09-08).} Take \(k=1\), \(N=r=2\),
nodes \(\zeta_1=1\), \(\zeta_2=-1\) and targets
\(Y_1=\operatorname{span}e_1\), \(Y_2=\operatorname{span}e_2\). A
degree-\(0\) curve is impossible (\(Y_1\ne Y_2\)) and
\(Q(z)=(z+1,\allowbreak{}\ z-1)^T\) has \(Q(1)=(2,\allowbreak{}0)^T\), \(Q(-1)=(0,\allowbreak{}-2)^T\), so
\(\delta_{\mathrm{Gr}}=d_{\min}=1\). The pencil
\(P(z)=(z+1,\allowbreak{}\ z(z-1))^T\) is admissible (\(P(1)=(2,\allowbreak{}0)^T\),
\(P(-1)=(0,\allowbreak{}2)^T\)) and its coordinates are coprime, so \(g_P=1\) and
\(P\) has no base point, but \(\operatorname{fdeg}(P)=2\). On
\(\mathbb T\), \(P^\sharp P=|z+1|^2+|z-1|^2=4\), so \(H_P=2\) and the
compiler returns \(P/2\), a polynomial column of McMillan degree
\(2=\operatorname{fdeg}(P)+0\), as the formula predicts, while the inner
matrix
\[S_1(z)=\frac12\begin{pmatrix}z+1&z-1\\ z-1&z+1\end{pmatrix},\qquad\det S_1=z,\]
is a router of degree \(1\) (\(S_1(1)e_1=e_1\), \(S_1(-1)e_1=-e_2\)),
and \(S_2:=\operatorname{diag}(1,\allowbreak{}z)S_1\) is inner of degree \(2\) with
first column \(P/2\) and \(\det S_2=z^2\). Here \(P\) is itself a
minimal basis of its column space \(V=\mathbb C(z)P\), so \(F_0=P/2\),
\(U_i=1\), and Theorem 5.7 correctly says that \(P/2\) is the
minimum-degree inner column \emph{with column space \(V\)}; the point is
that the column space of the optimal router's first column,
\(\mathbb C(z)Q\), is a different line. Unitarity of \(S_1,\allowbreak{}S_2\) on
\(\mathbb T\), the determinants, the gcd, the interpolation values and
the three McMillan degrees (Hankel ranks \(1,\allowbreak{}2,\allowbreak{}2\)) were verified
symbolically ({\small\texttt{repro/\allowbreak{}scratch\_\allowbreak{}v2/\allowbreak{}taller\_\allowbreak{}and\_\allowbreak{}graph\_\allowbreak{}witness.\allowbreak{}py}};
the workshop's own script {\small\texttt{check\_\allowbreak{}core.\allowbreak{}py}} is rerun in
{\small\texttt{repro/\allowbreak{}taller\_\allowbreak{}check/\allowbreak{}}}).

Section 8 verifies the degree formula on nine witnesses, all with
\(\operatorname{fdeg}(P)=\delta_{\mathrm{Gr}}=3\): one, two and double
base points in \(\mathbb D\) (cost \(1,\allowbreak{}2,\allowbreak{}2\)), the pair \(\pm\tfrac12\)
(\(2\)), one inside and one outside (\(1\)), a non-diagonal \(U_P\) with
\(\det U_P=z-\tfrac12\) (\(1\)), a unimodular \(U_P\) raising the column
degrees but not the Forney degree (\(0\)), and base points at \(2\) and
at \(\infty\) (\(0\)). Base points on \(\mathbb T\) were not tested
numerically (the Wilson iteration used there requires
\(P^\sharp P\succ0\) on \(\mathbb T\)).

\subsection{7. Consequences and comparison with the published
laws}\label{consequences-and-comparison-with-the-published-laws}

\textbf{7.1 Lines (\(k=1\)).} \(\mathrm{Gr}(1,\allowbreak{}H)=\mathbb P(H)\), minors
are coordinates, \(\operatorname{fdeg}\) is the projective degree, and
Theorem 1 is 6M3V Theorem 2.3 with a lower-bound proof that does not use
\(d\le L-1\).

\textbf{7.2 Direct-sum words.} Let \(w\) be a word over targets \(Y_a\)
in direct-sum position and \(F_0:=\sum_ap_a(z)W_a\) with
\(p_a:=\prod_{i:w_i\ne a}(z-\zeta_i)\) (the compiler input of 52B6,
Section 3). Every column has degree \(\le L-n_*\) and
\(F_0(\zeta_i)=p_{w_i}(\zeta_i)W_{w_i}\) has rank \(k\), so Lemma A
gives \(\delta_{\mathrm{Gr}}\le k(L-n_*)\). The 52B6 lower bound
transfers to curves: with \(D_a\) the block-dual detectors
(\(D_aW_b=\delta_{ab}I_k\)) and \(M\) a minimal basis through the data,
\(\det(D_aM(z))=\sum_I\det(D_a)_{:,\allowbreak{}I}\operatorname{minor}_I(M)\) has
degree \(\le\sum_je_j\), vanishes to order \(\ge k\) at every node not
labelled \(a\) and is nonzero at nodes labelled \(a\), so
\(\sum_je_j\ge k(L-n_a)\). Hence \(\delta_{\mathrm{Gr}}=k(L-n_*)\) and
Theorem 1 recovers 52B6 Theorem 1.2. Theorem 1 also explains why 52B6
Lemma 3.1 bounded the compiler by \(k\cdot(L-n_*)\) only: the exact
count is \(\sum_je_j\), which equals \(k\max_je_j\) precisely for this
word construction.

\textbf{7.3 Generic \(k(L-1)\) and the span bound.} When \(r=kL\) each
node is its own letter, \(n_*=1\), and \(\delta_{\mathrm{Gr}}=k(L-1)\):
33BE Corollary 1.2 and 1WXC Theorem 3.1. In general the
\(\sum_j(e_j+1)=\delta_{\mathrm{Gr}}+k\) coefficient vectors of a
minimal basis span a space containing every \(M(\zeta_i)\), hence \(H\);
so \(\delta_{\mathrm{Gr}}\ge r-k\), which is the lower half of 33BE
Theorem 1.1 and 52B6 Theorem 5.1. For \(r<kL\) the generic value can be
strictly smaller than \(k(L-1)\): in \(\mathrm{Gr}(2,\allowbreak{}4)\) it is
\(2\lfloor L/2\rfloor<2(L-1)\) for every \(L\ge3\) (Theorem 7.4).

\textbf{Proposition 7.3 (generic lower bound).} \emph{Let
\(e_0:=\lfloor L(r-k)/r\rfloor\), the least \(e\) with
\(r(e+1)>L(r-k)\), and \(m_0:=r(e_0+1)-L(r-k)\in\{1,\allowbreak{}\dots,\allowbreak{}r\}\). Let
\(\Omega_1\subseteq\mathrm{Gr}(k,\allowbreak{}r)^L\) be the set of data with
\(V_e=0\) for all \(e<e_0\) and \(\dim V_{e_0}=m_0\). Then \(\Omega_1\)
is nonempty and Zariski-open, and at every point of \(\Omega_1\)}
\[\delta_{\mathrm{Gr}}\ \ge\ k\,e_0+\max(0,\,k-m_0).\]

\emph{Proof.} \(V_e\) is the kernel of the incidence map
\(\iota_e:\mathbb C[z]_{\le e}\otimes H\to\bigoplus_iH/Y_i\), of
dimensions \(r(e+1)\) and \(L(r-k)\); injectivity of \(\iota_e\) for
\(e<e_0\) and surjectivity of \(\iota_{e_0}\) are Zariski-open
conditions on the data. A witness for both:
\(Y_i=\operatorname{span}\{e_a:a\notin I_i\}\) with the \(r-k\) missed
coordinates \(I_i\) assigned cyclically, so that coordinate \(a\) is
missed by \(|I_a|\in\{\lfloor L(r-k)/r\rfloor,\allowbreak{}\lceil L(r-k)/r\rceil\}\)
nodes, \(\sum_a|I_a|=L(r-k)\). Then \(f\in V_e\) iff each coordinate
\(f_a\) (degree \(\le e\)) vanishes at the \(|I_a|\) distinct nodes
missing \(a\), so \(\dim V_e=\sum_a\max(0,\allowbreak{}e+1-|I_a|)\), which is \(0\)
for \(e+1\le e_0\le\min_a|I_a|\), and equals
\(\#\{a:|I_a|=e_0\}=r(e_0+1)-L(r-k)=m_0\) for \(e=e_0\) (if \(L(r-k)/r\)
is an integer all \(|I_a|=e_0\) and \(\dim V_{e_0}=r=m_0\)). On the open
set, every column of a minimal basis attaining \(\delta_{\mathrm{Gr}}\)
is a nonzero element of \(V_{e_j}\), so \(e_j\ge e_0\); the columns with
\(e_j=e_0\) are linearly independent over \(\mathbb C\) and lie in
\(V_{e_0}\), so there are at most \(m_0\) of them and the remaining
\(\max(0,\allowbreak{}k-m_0)\) columns have degree \(\ge e_0+1\). \(\square\)

For \(k=1\) this is the one-band generic value
\(\lfloor L(r-1)/r\rfloor\) of 6M3V Theorem 3.2, where it is also an
upper bound; for \(r=kL\) it is \(k(L-1)\).

\textbf{Theorem 7.4 (generic value for \(r=2k\); in particular
\(\mathrm{Gr}(2,\allowbreak{}4)\)).} \emph{For data in the ambient Grassmannian
\(\mathrm{Gr}(k,\allowbreak{}2k)\), let \(L\ge1\) and \(e:=\lfloor L/2\rfloor\). Then
\(\delta_{\mathrm{Gr}}=k\,\allowbreak{}e\) on a nonempty Zariski-open subset of
\(\mathrm{Gr}(k,\allowbreak{}2k)^L\). In particular
\(\delta_{\mathrm{Gr}}=2\lfloor L/2\rfloor\) generically on
\(\mathrm{Gr}(2,\allowbreak{}4)^L\), for every \(L\).}

\emph{Proof.} For \(L=1\), a constant curve through the single target
has degree \(0=ke\); its actual span has dimension \(k\). For \(L\ge2\),
work in the ambient space of dimension \(r=2k\); Remark 3.1 permits this
even for special data of smaller span. Here
\(e_0=\lfloor L(r-k)/r\rfloor=\lfloor L/2\rfloor=e\) and
\(m_0=2k(e+1)-kL\), which is \(2k\) for \(L=2e\) and \(k\) for
\(L=2e+1\); in both cases \(m_0\ge k\), so Proposition 7.3 gives
\(\delta_{\mathrm{Gr}}\ge ke\) at every point of \(\Omega_1\).

\emph{Explicit data.} Write vectors of \(\mathbb C^{2k}\) as \((u;\allowbreak{}v)\)
with \(u,\allowbreak{}v\in\mathbb C^k\), and for the given nodes put
\(Y_i:=\{(x;\allowbreak{}\zeta_i^{e}x):x\in\mathbb C^k\}\), the graph of
\(\zeta_i^eI_k\). These are the values at the nodes of the curve
\(z\mapsto\operatorname{graph}(z^eI_k)\), whose minimal basis
\(P_e:=\binom{I_k}{z^eI_k}\) has column degrees \((e,\allowbreak{}\dots,\allowbreak{}e)\) (rank
\(k\) at every \(z\); leading-coefficient matrix \(\binom{0}{I_k}\) when
\(e>0\), and \(\binom{I_k}{I_k}\) when \(e=0\), in both cases of rank
\(k\)) and \(k\times k\) minors that include \(1\) and \(z^{ke}\); so
\(\operatorname{fdeg}(P_e)=ke\) and \(\delta_{\mathrm{Gr}}\le ke\) for
these data.

\emph{The graph data lie in \(\Omega_1\).} A vector
\(f=(u;\allowbreak{}v)\in\mathbb C[z]^{2k}\) of degree \(\le e'\) satisfies
\(f(\zeta_i)\in Y_i\) iff \(w:=v-z^eu\) vanishes at \(\zeta_i\), and
\(\deg w\le e+e'\). If \(e'\le e-1\) then \(\deg w\le2e-1<L\), so
\(w=0\), \(v=z^eu\), and \(\deg v\le e'<e\) forces \(u=0\):
\(V_{e'}=0\). If \(e'=e\) and \(L=2e+1\), again \(w=0\) and \(v=z^eu\)
with \(\deg v\le e\) forces \(u\) constant, so
\(V_e=\{(u;\allowbreak{}z^eu):u\in\mathbb C^k\}\) has dimension \(k=m_0\). If
\(e'=e\) and \(L=2e\), then \(w=c\,\allowbreak{}\pi\) with
\(\pi:=\prod_i(z-\zeta_i)=\sum_t\pi_tz^t\) and \(c\in\mathbb C^k\), and
\(v=z^eu+c\pi\) has degree \(\le e\) iff its coefficients of
\(z^{e+1},\allowbreak{}\dots,\allowbreak{}z^{2e}\) vanish, i.e.~\(u_j=-\pi_{e+j}\,\allowbreak{}c\) for
\(j=1,\allowbreak{}\dots,\allowbreak{}e\), with \(u_0\in\mathbb C^k\) free: \(\dim V_e=2k=m_0\).
Hence the graph data are in \(\Omega_1\), and
\(\delta_{\mathrm{Gr}}=ke\) there.

\emph{Openness.} Let \(\Omega_2\supseteq\Omega_1\) be the open set where
\(\iota_e\) is surjective; over \(\Omega_2\) the kernel \(V_e\) is a
vector bundle of rank \(m_0\) (kernel of a surjective map between vector
bundles, the target being \(\bigoplus_i\mathcal Q_i\) with
\(\mathcal Q\) the universal quotient bundle). For each node \(i\) the
\(k\)-linear bundle map \(V_e^{\times k}\to\Lambda^k\mathbb C^{2k}\),
\((f_1,\allowbreak{}\dots,\allowbreak{}f_k)\mapsto f_1(\zeta_i)\wedge\dots\wedge f_k(\zeta_i)\),
vanishes identically on a closed subset, so the set
\(W\subseteq\Omega_2\) where none of the \(L\) maps vanishes identically
is open. On \(W\), choose for each \(i\) a nonzero coordinate \(q_i\) of
the \(i\)-th map; \(\prod_iq_i\) is a nonzero polynomial function on
\(V_e^{\times k}\), so some tuple has all wedges nonzero and Lemma A
(last form) gives \(\delta_{\mathrm{Gr}}\le ke\). The graph data lie in
\(\Omega_1\cap W\) (the columns of \(P_e\) are in \(V_e\) with
independent values everywhere), so \(\Omega_1\cap W\) is a nonempty open
set on which \(\delta_{\mathrm{Gr}}=ke\). \(\square\)

The dimensions \(\dim V_{e'}\) of the graph data (\(0\) for \(e'<e\),
\(m_0\) for \(e'=e\)) were also checked by exact linear algebra for
\(k=2,\allowbreak{}3\) and \(1\le L\le12\)
({\small\texttt{repro/\allowbreak{}scratch\_\allowbreak{}v2/\allowbreak{}taller\_\allowbreak{}and\_\allowbreak{}graph\_\allowbreak{}witness.\allowbreak{}py}}). Version 1
stated the theorem for \((2,\allowbreak{}4)\) and \(3\le L\le7\) only, with the upper
bound taken from the exact random witnesses of Section 8; those
witnesses (now in {\small\texttt{repro/\allowbreak{}witnesses/\allowbreak{}}}) remain an independent
confirmation that \emph{random} data over \(\mathbb Q(i)\) lie in the
open set of the theorem, and the exact values \(3,\allowbreak{}6,\allowbreak{}6,\allowbreak{}9\) for random
data with \((k,\allowbreak{}r)=(3,\allowbreak{}6)\), \(L=3,\allowbreak{}4,\allowbreak{}5,\allowbreak{}6\) (Table 2) confirm the case
\(k=3\).

The naive count \(\dim\{\text{degree-}e\text{ maps}\}=4+4e\ge4L\)
predicts \(L-1\) (e.g.~\(3\) for \(L=4\), \(5\) for \(L=6\)) and is
wrong: a cubic through four generic points would need splitting type
\((0,\allowbreak{}3)\) or \((1,\allowbreak{}2)\), i.e.~a nonzero element of \(V_1\), which is
generically \(0\). The possible relation to quantum Schubert calculus
{[}Be{]} is not used in this proof; the degree values here follow from
the interpolation argument above.

\emph{Comparison with fixed-domain curve counts.} Buch--Pandharipande
{[}BP-Tev, Section 1.1, Definition 1.1 and Theorem 1.4{]} study virtual
Tevelev degrees with the marked domain fixed. Under
\(\mathrm{Gr}(2,\allowbreak{}4)\cong Q^4\), genus zero and \(L\ge3\), the dimension
constraint is \(d=L-1\) and their quadric formula gives
\((1+(-1)^d)/2\): one for odd \(L\), zero for even \(L\). This parity
comparison helps place Theorem 7.4 in the existing literature. A
vanishing virtual count is not by itself nonexistence of an
interpolating morphism, nor does a nonzero virtual count prove the
minimum or existence for every fixed node configuration considered here.
Our minimum and fixed-node assertion follow from the explicit
interpolation proof above; the construction of a rational inner router
is an additional analytic step. No priority or equivalence claim is
inferred from this comparison.

\textbf{Conjecture 7.5 (generic law).} Generic
\(\delta_{\mathrm{Gr}}=k\,\allowbreak{}e_0+\max(0,\allowbreak{}k-m_0)\), i.e.~the bound of
Proposition 7.3 is attained. The upper bound requires a witness tuple
with \(\min(k,\allowbreak{}m_0)\) columns in \(V_{e_0}\) and the rest in
\(V_{e_0+1}\) having independent node values, a degree-\(k\) multilinear
rather than linear condition (quadratic when \(k=2\)); it is proved for
\(r=2k\) and every \(L\) (Theorem 7.4: there \(m_0\ge k\) and the
formula reads \(k\lfloor L/2\rfloor\)), certified by exact computation
for \((2,\allowbreak{}3,\allowbreak{}L=4,\allowbreak{}5,\allowbreak{}6)\) (values \(2,\allowbreak{}3,\allowbreak{}4\)), and confirmed redundantly on
random data for \((2,\allowbreak{}4,\allowbreak{}L\le7)\) and \((3,\allowbreak{}6,\allowbreak{}L=3,\allowbreak{}4,\allowbreak{}5,\allowbreak{}6)\) (values
\(2\lfloor L/2\rfloor\) and \(3,\allowbreak{}6,\allowbreak{}6,\allowbreak{}9\)). For \(k=1\) and for \(r=kL\)
the conjecture is the published law.

\textbf{7.6 Detector ranks cannot determine \(d_{\min}\).} 52B6 Theorem
5.2 bounds \(d_{\min}\ge\max\{r-k,\allowbreak{}\Delta\}\) with
\(\Delta=\max_D\sum_i(k-\operatorname{rank}(DW_i))\) over constant
\(k\times N\) detectors \(D\) with \(\operatorname{rank}(DW_i)=k\) for
some \(i\). For \(k=2\), \(N=4\) and generic data, \(\Delta\) is a
Schubert count: \(\operatorname{rank}(DW_i)<2\) means that \(\ker D\), a
point of \(\mathrm{Gr}(2,\allowbreak{}4)\), meets \(Y_i\) nontrivially (the divisor
class \(\sigma_1\)), and \(\operatorname{rank}(DW_i)=0\) means
\(\ker D=Y_i\), i.e.~\(D\) is the annihilator of \(Y_i\), which has rank
\(2\) on the other generic planes. Since \(\sigma_1^4=2\) and
\(\sigma_1^5=0\) in \(H^*(\mathrm{Gr}(2,\allowbreak{}4))\), for \(L=4\) an admissible
\(\ker D\) can meet at most three of the planes (the two planes meeting
all four give inadmissible \(D\)), and for \(L\ge5\) at most four. Hence
generically \(\Delta=3\) for \(L=4\) and \(\Delta=4\) for \(L=5,\allowbreak{}6\),
whereas \(\delta_{\mathrm{Gr}}=4,\allowbreak{}4,\allowbreak{}6\). So the hierarchy bound is tight
for \(L=5\) and strictly below \(d_{\min}\) for \(L=4\) and \(L=6\).
Moreover the structured data of Section 8 (four planes on a
\((1,\allowbreak{}2)\)-curve) also have \(r=4\) and \(\Delta=3\) but \(d_{\min}=3\):
the pair \((r,\allowbreak{}\Delta)\) takes the same value on data with \(d_{\min}=3\)
and with \(d_{\min}=4\). This answers 52B6 Limitation 8.2 negatively for
these statistics: off the direct-sum locus, \(d_{\min}\) is the degree
of a curve and is not a function of \(r\) and \(\Delta\). The numerical
detector search of Section 8 finds the witnesses \((1,\allowbreak{}1,\allowbreak{}1,\allowbreak{}2)\),
\((1,\allowbreak{}1,\allowbreak{}1,\allowbreak{}1,\allowbreak{}2)\) and \((1,\allowbreak{}1,\allowbreak{}1,\allowbreak{}1,\allowbreak{}2,\allowbreak{}2)\) predicted by this count.

\subsection{8. Experiments}\label{experiments}

All computations are in {\small\texttt{repro/\allowbreak{}}} (Section 10). Nodes are
Gaussian-rational points of \(\mathbb T\)
(\(\tfrac{3+4i}5,\allowbreak{}\tfrac{5+12i}{13},\allowbreak{}\tfrac{8+15i}{17},\allowbreak{}\dots\)); targets
are random Gaussian-integer frames of rank \(k\); exact arithmetic over
\(\mathbb Q(i)\) uses SymPy's {\small\texttt{DomainMatrix}}; numerics use
NumPy/SciPy in double precision.

\textbf{Table 1. Exact \(\delta_{\mathrm{Gr}}\) (author's script
{\small\texttt{grassmann\_\allowbreak{}degree\_\allowbreak{}exact.\allowbreak{}py}}, \(k=2\)).} ``Witness'' is a random
element of \(V_{e_1}\times V_{e_2}\) with its Forney degree and the gcd
of its minors. The exact nodes, frames, witness pencils and lower-bound
certificates of every row except the last are the files
{\small\texttt{repro/\allowbreak{}witnesses/\allowbreak{}*.\allowbreak{}json}} described after Table 3.

{\def\LTcaptype{none} % do not increment counter
\begin{longtable}[]{@{}
  >{\raggedright\arraybackslash}p{(\linewidth - 12\tabcolsep) * \real{0.14285714}}
  >{\raggedright\arraybackslash}p{(\linewidth - 12\tabcolsep) * \real{0.14285714}}
  >{\raggedright\arraybackslash}p{(\linewidth - 12\tabcolsep) * \real{0.14285714}}
  >{\raggedright\arraybackslash}p{(\linewidth - 12\tabcolsep) * \real{0.14285714}}
  >{\raggedright\arraybackslash}p{(\linewidth - 12\tabcolsep) * \real{0.14285714}}
  >{\raggedright\arraybackslash}p{(\linewidth - 12\tabcolsep) * \real{0.14285714}}
  >{\raggedright\arraybackslash}p{(\linewidth - 12\tabcolsep) * \real{0.14285714}}@{}}
\toprule\noalign{}
\begin{minipage}[b]{\linewidth}\raggedright
data
\end{minipage} & \begin{minipage}[b]{\linewidth}\raggedright
\(r\)
\end{minipage} & \begin{minipage}[b]{\linewidth}\raggedright
\(r-k\)
\end{minipage} & \begin{minipage}[b]{\linewidth}\raggedright
\(k(L-1)\)
\end{minipage} & \begin{minipage}[b]{\linewidth}\raggedright
\(\delta_{\mathrm{Gr}}\)
\end{minipage} & \begin{minipage}[b]{\linewidth}\raggedright
column degrees
\end{minipage} & \begin{minipage}[b]{\linewidth}\raggedright
witness fdeg / gcd
\end{minipage} \\
\midrule\noalign{}
\endhead
\bottomrule\noalign{}
\endlastfoot
random seeds 1, 2, 3; \(N=4\), \(L=4\) & 4 & 2 & 6 & \textbf{4} & (2,2)
& 4 / 1 \\
seed 1 embedded in \(\mathbb C^5\)
(\(H=\mathbb C^4\subsetneq\mathbb C^5\)) & 4 & 2 & 6 & \textbf{4} &
(2,2) & 4 / 1 \\
random \((1,\allowbreak{}2)\)-curve data, \(L=4\) & 4 & 2 & 6 & \textbf{3} & (1,2) &
3 / 1 \\
same construction, \(L=5\) & 4 & 2 & 8 & \textbf{3} & (1,2) & 3 / 1 \\
generic \(L=3\) & 4 & 2 & 4 & \textbf{2} & (1,1) & 2 / 1 \\
generic \(L=5\) & 4 & 2 & 8 & \textbf{4} & (2,2) & 4 / 1 \\
generic \(L=6\) & 4 & 2 & 10 & \textbf{6} & (3,3) & 6 / 1 \\
generic \(L=7\) & 4 & 2 & 12 & \textbf{6} & (3,3) & 6 / 1 \\
\(\mathrm{Gr}(2,\allowbreak{}3)\), generic \(L=4,\allowbreak{}5,\allowbreak{}6\) ({\small\texttt{extra\_\allowbreak{}checks.\allowbreak{}py}}) &
3 & 1 & 6, 8, 10 & 2, 3, 4 & (1,1), (1,2), (2,2) & --- \\
\end{longtable}
}

For every tuple with smaller sum the log records
\(\dim V_{e_1},\allowbreak{}\dim V_{e_2}\) and the failing wedge test; for generic
\(L=4\), \(\dim V_0=\dim V_1=0\) and \(\dim V_2=4\), so \((0,\allowbreak{}3)\) and
\((1,\allowbreak{}2)\) are impossible and \((2,\allowbreak{}2)\) works. The
\(\mathrm{Gr}(2,\allowbreak{}3)\cong\mathbb P^2\) values coincide with the exact
\(k=1\) projective degrees of the annihilator lines computed separately
and with \(\lfloor2L/3\rfloor\).

\textbf{Table 2. Reviewer's independent exact computation
({\small\texttt{review\_\allowbreak{}exact\_\allowbreak{}delta.\allowbreak{}py}}, own code, different nodes and
seeds).}

{\def\LTcaptype{none} % do not increment counter
\begin{longtable}[]{@{}
  >{\raggedright\arraybackslash}p{(\linewidth - 4\tabcolsep) * \real{0.3333}}
  >{\raggedright\arraybackslash}p{(\linewidth - 4\tabcolsep) * \real{0.3333}}
  >{\raggedright\arraybackslash}p{(\linewidth - 4\tabcolsep) * \real{0.3333}}@{}}
\toprule\noalign{}
\begin{minipage}[b]{\linewidth}\raggedright
configuration
\end{minipage} & \begin{minipage}[b]{\linewidth}\raggedright
published value
\end{minipage} & \begin{minipage}[b]{\linewidth}\raggedright
\(\delta_{\mathrm{Gr}}\) (exact)
\end{minipage} \\
\midrule\noalign{}
\endhead
\bottomrule\noalign{}
\endlastfoot
\(k=1\), three distinct coplanar lines & 1 (33BE Thm 7.2) & 1 \\
\(k=1\), lines \(y_1,\allowbreak{}y_1,\allowbreak{}y_2\) & 2 (33BE Thm 7.2) & 2 \\
\(k=1\), four lines, pattern \(3+1\) & 3 (6M3V Thm 11.1) & 3 \\
\(k=1\), binary word \(y_1y_2y_1y_2y_1\) & \(\max(n_0,\allowbreak{}n_\infty)=3\)
(6M3V Thm 9.2) & 3 \\
\(k=1\), generic lines in \(\mathbb C^3\), \(L=4,\allowbreak{}5,\allowbreak{}6,\allowbreak{}7\) &
\(\lfloor2L/3\rfloor=2,\allowbreak{}3,\allowbreak{}4,\allowbreak{}4\) (6M3V Thm 3.2) & 2, 3, 4, 4 \\
\(k=2\) in \(\mathbb C^3\), generic \(L=4,\allowbreak{}5,\allowbreak{}6\) & \(\lfloor2L/3\rfloor\)
by duality & 2, 3, 4 \\
\(k=2\), \(A\oplus B=\mathbb C^4\), words AAB / ABAB / AAAB / AABB /
AAABB / ABABA & \(k(L-n_*)=4/4/6/4/6/6\) (52B6 Thm 1.2) & 4 / 4 / 6 / 4
/ 6 / 6 \\
\(k=3\), \(N=6\), \(L=2\) generic (direct sum) & \(k(L-1)=3\) (33BE Cor
1.2) & 3 \\
\(k=3\), \(N=6\), \(L=3\) generic & Prop. 7.3: \(\ge3\); Conj. 7.5: 3 &
3 \\
\(k=2\), \(N=4\), generic \(L=3,\allowbreak{}4,\allowbreak{}5,\allowbreak{}6\) (fresh seeds) & sandwich
\(2\le d\le2(L-1)\) & 2, 4, 4, 6 \\
\(k=3\), \(N=6\), generic \(L=4,\allowbreak{}5,\allowbreak{}6\) (version 2, reviewer's code,
{\small\texttt{scratch\_\allowbreak{}v2/\allowbreak{}k3r6\_\allowbreak{}generic.\allowbreak{}py}}) & Theorem 7.4:
\(3\lfloor L/2\rfloor=6,\allowbreak{}6,\allowbreak{}9\) & 6, 6, 9 \\
\end{longtable}
}

\textbf{Table 3. Inner interpolants and base-point degree increments
({\small\texttt{inner\_\allowbreak{}interpolant\_\allowbreak{}numeric.\allowbreak{}py}}, {\small\texttt{extra\_\allowbreak{}checks*.\allowbreak{}py}},
{\small\texttt{review\_\allowbreak{}basepoint\_\allowbreak{}additivity.\allowbreak{}py}}).} Pipeline: exact witness
\(M\) → \(R=M^\sharp M\) → outer factor \(H_o\) by Wilson's Newton
iteration {[}Wi{]} on a 4096-point grid (7--60 iterations;
\(|H_o^\sharp H_o-R|\le3.5\cdot10^{-13}\), Fourier tail beyond degree
\(m\) \(\le7\cdot10^{-13}\), all zeros of \(\det H_o\) outside
\(\overline{\mathbb D}\)) → Taylor coefficients of \(F_0=MH_o^{-1}\)
(400 terms) → block-Hankel SVD (\(60\times60\) blocks; rank = McMillan
degree, with a clean gap, e.g.~singular values
\(0.921,\allowbreak{}0.82,\allowbreak{}0.603,\allowbreak{}0.446,\allowbreak{}0,\allowbreak{}0,\allowbreak{}0\) for seed 1) → ERA realization {[}JP{]}
→ observability-Gramian scaling (isometric realization) → null-space
augmentation to a unitary realization matrix → square inner \(S\) →
checks. ``res'' is the largest of \(|F_0^*F_0-I|_{\mathbb T}\),
\(|R^*R-I|\) (isometric realization), \(|S^*S-I|_{\mathbb T}\); the
subspace distances
\(\max_i\operatorname{dist}(S(\zeta_i)\operatorname{ran}X,\allowbreak{}Y_i)\) are
\(<10^{-12}\) in every row.

In Table 3, the two mixing matrices are
\(U_1=\begin{pmatrix}z-\frac12&1\\0&1\end{pmatrix}\) and
\(U_2=\begin{pmatrix}z-\frac12&z\\1&1\end{pmatrix}\).

{\def\LTcaptype{none} % do not increment counter
\begin{longtable}[]{@{}
  >{\raggedright\arraybackslash}p{(\linewidth - 12\tabcolsep) * \real{0.14285714}}
  >{\raggedright\arraybackslash}p{(\linewidth - 12\tabcolsep) * \real{0.14285714}}
  >{\raggedright\arraybackslash}p{(\linewidth - 12\tabcolsep) * \real{0.14285714}}
  >{\raggedright\arraybackslash}p{(\linewidth - 12\tabcolsep) * \real{0.14285714}}
  >{\raggedright\arraybackslash}p{(\linewidth - 12\tabcolsep) * \real{0.14285714}}
  >{\raggedright\arraybackslash}p{(\linewidth - 12\tabcolsep) * \real{0.14285714}}
  >{\raggedright\arraybackslash}p{(\linewidth - 12\tabcolsep) * \real{0.14285714}}@{}}
\toprule\noalign{}
\begin{minipage}[b]{\linewidth}\raggedright
witness
\end{minipage} & \begin{minipage}[b]{\linewidth}\raggedright
\(\delta_{\mathrm{Gr}}\)
\end{minipage} & \begin{minipage}[b]{\linewidth}\raggedright
predicted \(\deg\) (Thm 6.1)
\end{minipage} & \begin{minipage}[b]{\linewidth}\raggedright
Hankel rank of \(F_0\)
\end{minipage} & \begin{minipage}[b]{\linewidth}\raggedright
\(\operatorname{wind}\det S\)
\end{minipage} & \begin{minipage}[b]{\linewidth}\raggedright
Hankel rank of \(S\)
\end{minipage} & \begin{minipage}[b]{\linewidth}\raggedright
res
\end{minipage} \\
\midrule\noalign{}
\endhead
\bottomrule\noalign{}
\endlastfoot
random seed 1, minimal basis & 4 & 4 & 4 & 4.000000 & 4 &
\(1.2\cdot10^{-11}\) \\
random seed 2 & 4 & 4 & 4 & 4.000000 & 4 & \(7.5\cdot10^{-12}\) \\
seed 1 in \(\mathbb C^5\) & 4 & 4 & 4 & 4.000000 & 4 &
\(2.2\cdot10^{-12}\) \\
\((1,\allowbreak{}2)\)-curve data, minimal basis \([p_1,\allowbreak{}p_2]\) & 3 & 3 & 3 & 3.000000
& 3 & \(2.4\cdot10^{-15}\) \\
\([p_1,\allowbreak{}(z-\tfrac12)p_2]\) (base point in \(\mathbb D\)) & 3 & 4 & 4 &
4.000000 & 4 & \(5.7\cdot10^{-15}\) \\
\([p_1,\allowbreak{}(z-\tfrac12)(z-\tfrac13)p_2]\) & 3 & 5 & 5 & 5.000000 & 5 &
\(7.3\cdot10^{-15}\) \\
\((z-\tfrac12)[p_1,\allowbreak{}p_2]\) (double base point) & 3 & 5 & 5 & 5.000000 & 5
& \(8.0\cdot10^{-15}\) \\
\([(z-\tfrac12)p_1,\allowbreak{}(z-2)p_2]\) (one inside, one outside) & 3 & 4 & 4 &
4.000000 & 4 & \(3.4\cdot10^{-15}\) \\
\(P_0U_1\) (non-diagonal, \(\det U_1=z-\tfrac12\)) & 3 & 4 & 4 &
4.000000 & 4 & \(3.0\cdot10^{-15}\) \\
\(P_0U_2\) (\(\det U_2=-\tfrac12\), column degrees (2,3)) & 3 & 3 & 3 &
3.000000 & 3 & \(2.7\cdot10^{-15}\) \\
\([p_1,\allowbreak{}(z^2-\tfrac14)p_2]\) (base points \(\pm\tfrac12\)) & 3 & 5 & 5 &
5.000000 & 5 & \(6.4\cdot10^{-15}\) \\
\([p_1,\allowbreak{}(z-2)p_2]\) (base point outside \(\overline{\mathbb D}\)) & 3 & 3
& 3 & 3.000000 & 3 & \(2.2\cdot10^{-15}\) \\
\([p_1,\allowbreak{}z^2p_1+p_2]\) (base point at \(\infty\); leading matrix rank 1) &
3 & 3 & 3 & 3.000000 & 3 & \(2.2\cdot10^{-15}\) \\
generic \(L=5\) & 4 & 4 & 4 & 4.000000 & 4 & \(3.6\cdot10^{-11}\) \\
generic \(L=6\) & 6 & 6 & 6 & 6.000000 & 6 & \(5.0\cdot10^{-10}\) \\
\end{longtable}
}

The random instances have zeros of \(\det H_o\) as close as \(1.014\) to
\(\mathbb T\) (\(L=5\)), which is why their residuals are
\(10^{-10}\)--\(10^{-11}\) rather than \(10^{-15}\); the structured
instances (zeros at modulus \(\ge2\)) reach machine precision. In every
minimal-basis row \(\deg S=\delta_{\mathrm{Gr}}\) with all checks below
\(10^{-9}\), and every base-point row matches Theorem 6.1.

\textbf{Witness files and the computer-assisted parts.} Each file a JSON
file in {\small\texttt{repro/\allowbreak{}witnesses/\allowbreak{}}} contains, as exact SymPy strings over
\(\mathbb Q(i)\): the nodes, the \(N\times k\) frames of the targets,
the value \(\delta_{\mathrm{Gr}}\), the column degrees \((e_1,\allowbreak{}e_2)\),
the two witness columns \(f_1\in V_{e_1}\), \(f_2\in V_{e_2}\), and the
lower-bound certificate, i.e.~for every pair \(e_1\le e_2\) with
\(e_1+e_2<\delta_{\mathrm{Gr}}\) the dimensions
\(\dim V_{e_1},\allowbreak{}\dim V_{e_2}\) (a zero dimension rules the pair out) and,
when both are positive, the fact that some wedge form
\(V_{e_1}\times V_{e_2}\to\Lambda^2Y_i\) is the zero matrix. The script
{\small\texttt{verify\_\allowbreak{}witness.\allowbreak{}py}} recomputes everything from the file
(interpolation and rank at the nodes,
\(\operatorname{fdeg}=\delta_{\mathrm{Gr}}\) with coprime minors, the
dimensions, the zero wedge forms) in about 15 seconds; for a reader
without a computer, the \(L=3\) file is small enough to check by hand
(\(V_0=0\) because three generic planes have no common vector, and
\(\dim V_1=2\)). In version 2 no theorem depends on these files: Theorem
7.4 is proved with the explicit graph witness. The computer-assisted
statements are the values in Tables 1 and 2 for \emph{random} data
(exact certificates, in the sense of Section 10), the detector patterns
of Section 7.6 (exact for the annihilator witnesses, numerical for the
search), everything in Table 3 (floating point, with the residuals
shown), and the Potapov search (evidence only).

\textbf{Detector ranks (numerical, {\small\texttt{detector\_\allowbreak{}bound\_\allowbreak{}plucker}}).}
For \(k=2\), \(N=4\): generic \(L=4\) (seeds 1, 2): \(\Delta=3\) with
pattern \((1,\allowbreak{}1,\allowbreak{}1,\allowbreak{}2)\), \(\delta_{\mathrm{Gr}}=4\); \((1,\allowbreak{}2)\)-curve data
\(L=4\): \(\Delta=3\) \((1,\allowbreak{}1,\allowbreak{}1,\allowbreak{}2)\), \(\delta_{\mathrm{Gr}}=3\); generic
\(L=5\): \(\Delta=4\) \((1,\allowbreak{}1,\allowbreak{}1,\allowbreak{}1,\allowbreak{}2)\), \(\delta_{\mathrm{Gr}}=4\);
generic \(L=6\): \(\Delta=4\) \((1,\allowbreak{}1,\allowbreak{}1,\allowbreak{}1,\allowbreak{}2,\allowbreak{}2)\),
\(\delta_{\mathrm{Gr}}=6\). The search parametrises admissible detectors
by the Klein quadric (each rank condition is a hyperplane section, since
\(\det(DW_i)\) is linear in the Plücker vector of the row space of
\(D\)) with rank tolerances \(10^{-8}\); the values agree with the
Schubert count of Section 7.6.

\textbf{Direct search over Blaschke--Potapov products (reviewer,
{\small\texttt{review\_\allowbreak{}potapov\_\allowbreak{}search.\allowbreak{}py}}).} As evidence independent of the
proof, all rational inner \(4\times4\) functions of degree \(d\) were
parametrised as \(U\prod_{q=1}^d(I+(b_{a_q}-1)u_qu_q^*)\) and the
objective \(\sum_i\|(I-\Pi_{Y_i})S(\zeta_i)X\|_F^2\) minimised by
Levenberg--Marquardt from 120 random starts per degree. For the \(L=3\)
instance (\(\delta_{\mathrm{Gr}}=2\)): best values \(1.99\) (\(d=0\)),
\(0.967\) (\(d=1\)), \(6\cdot10^{-30}\) (\(d=2\)). For the \(L=4\)
instance (\(\delta_{\mathrm{Gr}}=4\), where the prior bounds \(r-k=2\)
and \(\Delta=3\) are not tight): \(0.572\) (\(d=2\)), \(0.150\)
(\(d=3\), 120 restarts), \(3\cdot10^{-30}\) (\(d=4\)). The search never
found a degree-\(3\) interpolant, and found one of degree \(4\), as
Theorem 1 requires.

\subsection{9. Open questions}\label{open-questions}

\begin{enumerate}
\def\labelenumi{\arabic{enumi}.}
\tightlist
\item
  \textbf{Generic upper bound.} Prove Conjecture 7.5 for \(r\ne2k\),
  i.e.~that for generic data a tuple in \(\prod_jV_{e_j}\) with
  \(e_j\in\{e_0,\allowbreak{}e_0+1\}\) has independent node values. The graph witness
  of Theorem 7.4 (nodes on the curve
  \(z\mapsto\operatorname{graph}(z^eI_k)\)) lies in \(\Omega_1\) only
  when \(r=2k\); for other \((k,\allowbreak{}r)\) one needs a curve of degree
  \(ke_0+\max(0,\allowbreak{}k-m_0)\) whose node values are in \(\Omega_1\). The
  quantum Schubert calculus of \(\mathrm{Gr}(k,\allowbreak{}r)\) {[}Be, RRW2{]} may
  help study this question. Such a comparison must retain the fixed
  marked nodes on the domain; no general identification with an ordinary
  Gromov-Witten invariant is asserted here.
\item
  \textbf{Singular strata.} Theorem 1 reduces \(d_{\min}\) on every
  stratum (overlapping bands, repeated targets, incidence tables) to the
  degree of a curve through Schubert conditions, computable exactly by
  the enumeration of Lemma A but without a closed formula; the block law
  of 6M3V Theorem 3.2 for orthogonal bands and the occupancy law of 52B6
  are the two closed cases known.
\item
  \textbf{Delay and error at fixed degree for \(k>1\).} The rank--error
  dichotomy and the resonant trichotomy of 6M3V (Theorems 4.1 and 5.1)
  should generalise with the wedge condition
  \(f_1(\zeta_i)\wedge\dots\wedge f_k(\zeta_i)\ne0\) replacing
  primitivity: \(E_d=0\) iff some tuple in \(\prod_jV_{e_j}\) with
  \(\sum e_j\le d\) has independent node values \emph{except} at a set
  of nodes that can be deleted as base points. The border degree
  \(d_\partial\) and the divergence of the peak Wigner--Smith delay for
  \(d_\partial\le d<d_{\min}\) are not worked out here.
\item
  \textbf{Real and orthogonal versions.} For real data and
  real-orthogonal routers, Sottile's real Schubert calculus {[}So{]}
  suggests that the real \(\delta_{\mathrm{Gr}}\) can exceed the complex
  one; the Fejér--Riesz and completion steps preserve realness, so the
  real law should again be the degree of a real curve.
\item
  \textbf{Confluent nodes and fixed endpoint frames} are not addressed.
\end{enumerate}

\subsection{10. Reproducibility}\label{reproducibility}

The directory {\small\texttt{repro/\allowbreak{}}} accompanying this manuscript contains the
author's scripts ({\small\texttt{scratch\_\allowbreak{}D2/\allowbreak{}}}:
{\small\texttt{grassmann\_\allowbreak{}degree\_\allowbreak{}exact.\allowbreak{}py}},
{\small\texttt{inner\_\allowbreak{}interpolant\_\allowbreak{}numeric.\allowbreak{}py}}, {\small\texttt{extra\_\allowbreak{}checks.\allowbreak{}py}},
{\small\texttt{extra\_\allowbreak{}checks2.\allowbreak{}py}}), the reviewer's scripts
({\small\texttt{scratch\_\allowbreak{}D2\_\allowbreak{}review/\allowbreak{}}}: {\small\texttt{review\_\allowbreak{}exact\_\allowbreak{}delta.\allowbreak{}py}},
{\small\texttt{review\_\allowbreak{}basepoint\_\allowbreak{}additivity.\allowbreak{}py}},
{\small\texttt{review\_\allowbreak{}potapov\_\allowbreak{}search.\allowbreak{}py}}), the exact witness files with their
verifier ({\small\texttt{witnesses/\allowbreak{}*.\allowbreak{}json}}, {\small\texttt{verify\_\allowbreak{}witness.\allowbreak{}py}},
{\small\texttt{export\_\allowbreak{}witnesses.\allowbreak{}py}}), the version-2 checks
({\small\texttt{scratch\_\allowbreak{}v2/\allowbreak{}}}: symbolic verification of Example 6.2 and of the
graph witness of Theorem 7.4, the \((3,\allowbreak{}6)\) instances, the rerun of the
\(L=3\) Potapov stage), the workshop's script and its rerun
({\small\texttt{taller\_\allowbreak{}check/\allowbreak{}}}), the original logs and the logs of every
rerun, with commands, runtimes and expected final lines in
{\small\texttt{repro/\allowbreak{}README.\allowbreak{}md}}. Environment: Python 3.12.6, SymPy 1.14.0,
NumPy 2.5.1, SciPy 1.18.0. The exact computations are certificates: a
witness tuple with nonzero wedges at all nodes is an exact upper bound,
and the failure of every tuple of smaller degree sum is an exact lower
bound. The numerical pipeline is a check of the two classical inputs
(Fejér--Riesz factor, isometric realization and completion) on concrete
instances, not part of the proof.

\subsection{AI-assistance statement}\label{ai-assistance-statement}

The received version declares substantial generative-AI assistance in
drafting, mathematical derivations, code and review, attributed there to
``Claude Fable 5.1'' (Anthropic). That historical model identity, the
reported review sessions and their independence have not been
authenticated by this workshop. Historical descriptions of an ``author''
or ``reviewer'' script identify the supplied files, not a verified
independent editorial review. OpenAI Codex assisted with the present
private workshop revision: checking arguments and sources, correcting
identified errors, running the checks recorded in the accompanying
revision report, and preparing the PDF. Numerical checks support only
their tested instances. This assistance is not an independent human peer
review or a formal AIRR model-assessment report. Lluis Eriksson is the
declared author and AIRR founder/operator; this conflict must be handled
by the separate editorial process. Scientific authorship and approval
remain with the author.

The deposited v003 was subsequently assessed in actual OpenAI Codex
calls by gpt-6-astra and gpt-5.6-sol on 9 September 2026. Their original
reports and subsequent clarifications are retained. Both sets of
comments informed this v004 revision; Astra had also participated in
earlier manuscript preparation. A separate reviewing context does not
erase that history. Neither the v003 reports nor this disclosure
constitute approval of the revised v004 artifact.

\subsection{Changes in internal revision v004 (9 September
2026)}\label{changes-in-internal-revision-v004-9-september-2026}

\begin{enumerate}
\def\labelenumi{\arabic{enumi}.}
\tightlist
\item
  Explained cancellation of removable boundary singularities using
  boundedness on a punctured unit-circle arc.
\item
  Added a fixed-version comparison with Buch--Pandharipande's virtual
  Tevelev degrees, separating parity context from existence, minimality
  and the analytic router construction.
\end{enumerate}

Made the already cited AJL completion theorem more readily checkable by
adding its page, subparts, equations and the dimension correspondence;
retained the existing inversion-of-variable argument.

The earlier deposited v003 and its model reports are retained. This
revision responds to their findings; it is not a final acceptance.

\subsection{Changes in internal revision v003 (8 September
2026)}\label{changes-in-internal-revision-v003-8-september-2026}

\begin{enumerate}
\def\labelenumi{\arabic{enumi}.}
\tightlist
\item
  Corrected matrix multiplication order in Lemma 5.3.
\item
  Made the degree identity for an arbitrary minimal basis explicit
  before Theorem 5.7, so that the additivity proof does not rely on
  interpolation optimality.
\item
  Separated the \(L=1\) constant-curve case, its actual data span, and
  the degree-zero leading coefficient in Theorem 7.4.
\item
  Corrected the degree of the multilinear witness condition in
  Conjecture 7.5 and clarified the base-point discussion at infinity.
\item
  Qualified historical AI-review provenance; the new numerical and exact
  checks are separately recorded in the private revision report.
\end{enumerate}

\subsection{Changes in version 2 (8 September
2026)}\label{changes-in-version-2-8-september-2026}

Answers to the AIRR workshop evaluation of {\small\texttt{paper3.\allowbreak{}pdf}}
(TALLER-0003, 2026-09-08); the point-by-point response is
{\small\texttt{taller\_\allowbreak{}feedback/\allowbreak{}RESPUESTA-TALLER-0003.\allowbreak{}md}}.

\begin{enumerate}
\def\labelenumi{\arabic{enumi}.}
\tightlist
\item
  \emph{Theorem 6.1, last clause (workshop: proven flaw; accepted).} The
  clause ``the compiler attains \(d_{\min}\) iff \(P\) has no base point
  in \(\mathbb D\)'' and the last sentence of the proof, which used
  \(\operatorname{fdeg}(P)=\delta_{\mathrm{Gr}}\) for every admissible
  \(P\), were wrong: admissibility means interpolation, not minimality.
  Corrected to ``iff \(\operatorname{fdeg}(P)=\delta_{\mathrm{Gr}}\) and
  \(P\) has no base point in \(\mathbb D\)'', with the proof fixed; the
  degree formula
  \(\deg(PH_P^{-1})=\operatorname{fdeg}(P)+\#\{\text{base points in }\mathbb D\}\)
  and Theorem 1 are unaffected. The workshop's example (\(k=1\), nodes
  \(\pm1\), \(P=(z+1,\allowbreak{}z(z-1))^T\), routers \(S_1\), \(S_2\)) is
  reproduced and verified as Example 6.2. Abstract and Section 1.2 (item
  3) reworded accordingly.
\item
  \emph{Scope of the generic value \(2\lfloor L/2\rfloor\) (workshop:
  restrict or prove; accepted, proved).} Theorem 7.4 is now proved for
  all \(L\ge1\), and for \(\mathrm{Gr}(k,\allowbreak{}2k)\) with value
  \(k\lfloor L/2\rfloor\), by the explicit witness
  \(Y_i=\operatorname{graph}(\zeta_i^{\lfloor L/2\rfloor}I_k)\) and the
  lower bound of Proposition 7.3, whose statement now names the open set
  \(\Omega_1\) on which it holds. The abstract, Sections 1.2, 7.3,
  Conjecture 7.5 and Open question 1 were updated; the
  quantum-cohomology paragraph is kept as a consistency remark.
\item
  \emph{Witnesses and computer-assisted parts (workshop: evidence not
  attached; accepted).} The exact data, witness pencils and lower-bound
  certificates behind Table 1 are now data files
  {\small\texttt{repro/\allowbreak{}witnesses/\allowbreak{}*.\allowbreak{}json}} with an independent verifier,
  described in Section 8, which also states which claims are
  computer-assisted. Theorem 7.4 no longer depends on them. Table 2
  gains the \((3,\allowbreak{}6)\), \(L=4,\allowbreak{}5,\allowbreak{}6\) row.
\item
  \emph{Comparison with all-pass subspace interpolation (workshop:
  compare in more detail; accepted).} Section 1.1 now describes the
  boundary SNIP of Bharath et al.~(data, existence criterion, degree of
  the construction) and states the three differences with the routing
  problem; {[}ABKW{]} added for unconstrained minimal-degree
  interpolation.
\item
  \emph{Reproducibility.} {\small\texttt{repro/\allowbreak{}README.\allowbreak{}md}} rewritten with
  commands, runtimes and expected final lines; all fast scripts rerun on
  2026-09-08 (identical verdicts), the workshop's
  {\small\texttt{check\_\allowbreak{}core.\allowbreak{}py}} rerun, the \(L=3\) stage of the Potapov search
  rerun.
\item
  Date line and AI-assistance statement updated. No other mathematical
  statement changed.
\end{enumerate}

\subsection{References}\label{references}

{[}ABKW{]} A. C. Antoulas, J. A. Ball, J. Kang, J. C. Willems, \emph{On
the solution of the minimal rational interpolation problem}, Linear
Algebra Appl. 137/138 (1990), 511--573.

{[}AJL{]} D. Alpay, P. E. T. Jorgensen, I. Lewkowicz,
\emph{Characterizations of rectangular (para)-unitary rational
functions}, Opuscula Math. 36 (2016), no. 6, 695--716; arXiv:1410.0283.
Theorem 2.1.

{[}Be{]} A. Bertram, \emph{Quantum Schubert calculus}, Adv. Math. 128
(1997), 289--305.

{[}BGAP{]} A. S. Bharath, D. S. Gaharwar, K. Appaiah, D. Pal,
\emph{Design of discrete-time matrix all-pass filters using subspace
Nevanlinna--Pick interpolation}, Signal Processing 204 (2023), 108839.

{[}BD{]} V. Bolotnikov, H. Dym, \emph{On Boundary Interpolation for
Matrix Valued Schur Functions}, Mem. Amer. Math. Soc. 181 (2006), no.
856.

{[}DR{]} M. A. Dritschel, J. Rovnyak, \emph{The operator Fejér--Riesz
theorem}, in: A Glimpse at Hilbert Space Operators: Paul R. Halmos in
Memoriam, Oper. Theory Adv. Appl. 207, Birkhäuser, Basel, 2010,
223--254.

{[}Fo{]} G. D. Forney, Jr., \emph{Minimal bases of rational vector
spaces, with applications to multivariable linear systems}, SIAM J.
Control 13 (1975), 493--520.

{[}JP{]} J.-N. Juang, R. S. Pappa, \emph{An eigensystem realization
algorithm for modal parameter identification and model reduction}, J.
Guidance Control Dynam. 8 (1985), 620--627.

{[}Ka{]} T. Kailath, \emph{Linear Systems}, Prentice-Hall, 1980, §6.5.4.

{[}MH{]} C. Martin, R. Hermann, \emph{Applications of algebraic geometry
to systems theory: the McMillan degree and Kronecker indices of transfer
functions as topological and holomorphic system invariants}, SIAM J.
Control Optim. 16 (1978), 743--755.

{[}Po{]} V. P. Potapov, \emph{The multiplicative structure of
J-contractive matrix functions}, Trudy Moskov. Mat. Obshch. 4 (1955),
125--236; English transl. Amer. Math. Soc. Transl. (2) 15 (1960),
131--243.

{[}RRW1{]} M. S. Ravi, J. Rosenthal, X. Wang, \emph{Dynamic pole
assignment and Schubert calculus}, SIAM J. Control Optim. 34 (1996),
813--832.

{[}RRW2{]} M. S. Ravi, J. Rosenthal, X. Wang, \emph{Degree of the
generalized Plücker embedding of a Quot scheme and quantum cohomology},
Math. Ann. 311 (1998), 11--26.

{[}Ro{]} M. Rosenblum, \emph{Vectorial Toeplitz operators and the
Fejér--Riesz theorem}, J. Math. Anal. Appl. 23 (1968), 139--147.

{[}So{]} F. Sottile, \emph{Rational curves on Grassmannians: systems
theory, reality, and transversality}, in: Advances in Algebraic Geometry
Motivated by Physics, Contemp. Math. 276, Amer. Math. Soc., 2001, 9--42;
arXiv:math/0012079.

{[}Wi{]} G. T. Wilson, \emph{The factorization of matricial spectral
densities}, SIAM J. Appl. Math. 23 (1972), 420--426.

ARR-2026-07EH3CZRD89YC8JN, L. Eriksson, \emph{Lossless Calibration Is
Stored Memory: A Topological McMillan-Degree and Wigner--Smith Law},
2026 (Lemma 3.2).

ARR-2026-1WXCVX96S68GA9VK, L. Eriksson, \emph{Orthogonal Spectral
Fan-Out Costs \(k(L-1)\) States: a Global Memory Law Beyond Pairwise
Routing}, 2026 (Theorem 3.1, Lemma 3.2).

ARR-2026-33BE6K3HW78F4T7F, L. Eriksson, \emph{Exact Memory of Direct-Sum
Spectral Fan-Out: Generic Maximality, Colliding Targets, and the
Three-Line Phase Diagram}, 2026 (Theorem 1.1, Corollary 1.2, Theorem
7.2).

ARR-2026-52B6MSS1W197W9T2, L. Eriksson, \emph{Exact Memory of Finite
Spectral Routing Tables: The Direct-Sum Occupancy Law and Its Singular
Strata}, 2026 (Theorem 1.2, Lemma 2.1, Lemma 3.1, Theorems 5.1--5.2,
Limitation 8.2).

ARR-2026-6M3VGTXZ6W8JW9C9, L. Eriksson, \emph{Projective Memory and
Resonant Bottlenecks in Passive Spectral Routing}, 2026 (Definition 2.1,
Lemma 2.2, Theorem 2.3, Corollary 2.4, Theorems 3.2, 4.1, 5.1, 9.2,
11.1).

{[}BP-Tev{]} A. S. Buch and R. Pandharipande, \emph{Tevelev degrees in
Gromov-Witten theory}, arXiv:2112.14824v1 {[}math.AG{]}, 2021, Sections
1.1 and 1.3, Theorem 1.4.
\href{https://arxiv.org/abs/2112.14824v1}{Fixed version}.

\end{document}
